% Standalone, unreviewed research manuscript. Compile with XeLaTeX.
\documentclass[11pt,a4paper]{article}
\usepackage[margin=23mm]{geometry}
\usepackage{fontspec}
\IfFontExistsTF{Noto Sans CJK JP}{\setmainfont{Noto Sans CJK JP}}{%
\setmainfont{IPAexMincho}[BoldFont=IPAexMincho,ItalicFont=IPAexMincho,
BoldFeatures={FakeBold=2},ItalicFeatures={FakeSlant=.15}]}
\setmonofont{DejaVu Sans Mono}[Scale=.8]
\XeTeXlinebreaklocale "ja"
\XeTeXlinebreakskip=0pt plus 1pt
\usepackage{amsmath,amssymb,amsthm,mathtools,longtable,booktabs,array,enumitem,fancyvrb,xcolor,xurl}
\usepackage[colorlinks=true,linkcolor=blue!50!black,urlcolor=blue!50!black]{hyperref}
\setlength{\parindent}{0pt}
\setlength{\parskip}{.5em}
\setlength{\emergencystretch}{3em}
\setlist{nosep,leftmargin=1.7em}
\newtheorem{definition}{定義}[section]
\newtheorem{lemma}[definition]{補題}
\newtheorem{theorem}[definition]{定理}
\newtheorem{proposition}[definition]{命題}
\renewcommand{\proofname}{証明}
\renewcommand{\contentsname}{目次}
\newcommand{\Ret}{\mathsf{return}}
\newcommand{\Call}{\mathsf{call}}
\newcommand{\Delay}{\mathsf{delay}}
\newcommand{\Force}{\mathsf{force}}
\newcommand{\Scope}{\mathsf{scope}}
\newcommand{\Req}{\mathsf{Req}}
\newcommand{\Bad}{\mathsf{Bad}}
\newcommand{\Act}{\mathsf{Active}}
\newcommand{\El}{\mathsf{El}}
\newcommand{\fv}{\operatorname{fv}}
\newcommand{\dom}{\operatorname{dom}}
\newcommand{\Core}{\mathcal K}
\newcommand{\Src}{\mathcal S}
\newcommand{\D}{\mathcal D}
\title{Yulang のソース導出と実行核の接続\
\large コンストラクタ多相性、保護の由来、再開時の寿命と選択}
\author{単独構成・自己点検：このセッションの主エージェント}
\date{2026年10月24日改訂}
\begin{document}
\maketitle

\section*{到達点と訂正}

本稿は、推論器に依存しない型健全性を、ソースの宣言的な導出へ戻すための後続稿である。
\nolinkurl{yulang-operational-type-safety-ja.tex} の第12節までの体系を
\(\Core_{\rm old}\)、\nolinkurl{yulang-dot-selection-safety-ja.tex} の有限証拠構成を選択稿と呼ぶ。
基準 commit は \texttt{6bf6ec8a82eadaa75467e896f5a2658d0c3f12f8}。
旧操作稿の SHA-256 は
\nolinkurl{bba63cf35153335a22be221ad0925dcef7ac4afdf441e49a5fa5bb2a555f799c}、
旧選択稿は
\nolinkurl{66e5da71711a6457a9b77915c9484f5a3dc23659c1f96fdcdbf258240737a780}。
元のユーザー文書は改変しない。

今回、単なる未証明の橋だけでなく、旧機械と既存の source-owner 寿命の具体化との不一致を特定した。
保存した Boundary の receiver を、再開した code の owner と同時に改名すると、
元の receiver が既に終了していても保護が復活する。
旧稿の値型定理を反駁するものではないが、そのまま Yulang の衛生性へ適用する橋は成立しない。
第\ref{sec:owners}節で、過去の証拠を固定したまま実行 owner だけを解決する機械
\(\Core_{\rm life}\) を定義し、差分の保存・進行と有限表示を示す。

ソース側では、コンストラクタの多相 payload の導入・取得・各使用時の型適用を、
完成した核 code を前提としない宣言的規則から構成する。
一般の source と \texttt{'e?} の全対応を、実装だけの残務として隠さない。
特に、複数の独立な入力 occurrence と一つの印の対応を独立の義務として判定する。
\texttt{compose} の直接呼出しについては、第\ref{sec:named}節で既存の source frame 構成を適用し、
最後の formal \(x\) の実行 entry から scope/receiver を作る具体化まで進める。
この bounded な対応を alias/anonymous/一般の捕捉へ拡張する義務は残る。
今回の追補（第\ref{sec:capture}--\ref{sec:source-family}節）は、source の自由変数と
実際の値経路から捕捉を構成する。既存 evidence を複写せず、未実行 recipe の binder と分離する。
追補の基準は \texttt{c1ac52b0a} に保存した前版の source 稿である。
さらに non-Var 呼出しを扱う一つの全域な候補政策を明記し、候補に対する source 定理へ合成する。
これは旧 Oracle の全構文についての対応や、未選択の意味論の採用を意味しない。
独立査読・機械検証は未実施。以下の紙上導出は証明 gate の昇格ではない。

\tableofcontents

\section{ソースの導出に何を含めるか}
\label{sec:source}

\subsection{構文と原文への対応}

\(\Src\) は元稿の名前解決後の13式、11 pattern、3 statement の構文である。
parser が読む文字列とこの構文との全対応は、別に列挙する表層展開の義務である。
本稿の source は核の \(\Scope\)、\(\mathsf{release}\)、\(\mathsf{View}\)、
動的 receiver 名を持たない。名前は lexical binder/declaration identity に解決されている。
型注釈の \texttt{for 'a: T} は型の全称 binder であり、項に新しい type-lambda 記法を要求しない。
各注釈には元の source occurrence が残り、推論で選んだ型は明示注釈へ変わらない。

型の文法は操作稿の値型と \(\forall a.A\)、遅延型 \(U(E,A)\)、
四ポート Function、qualified telescope を用いる。
型の各 node は有限の構文または guarded な宣言 template の参照である。
nominal の型引数へ \(\forall\) を置く形も明示 instance として表せるが、
これは impredicative な宣言的 instance であり、その推論の完全性を意味しない。

\subsection{証明木と推論の区別}

宣言的判断を
\[
 \Delta;\Gamma;Q\vdash_s e:A\mid d,
 \qquad \Delta;\Gamma;Q\vdash_v n:A\mid d_v
\]
と書く。\(d,d_v\) は下記の有限の規則木であり、出力 code、scope 表、
実行安全性を premise に持たない。\(n\) は非拡張な source introduction、または既存の値 binder である。
非拡張 introduction は literal、lambda、明示 delay、非拡張な子を持つ tuple/record/variant、
これらの型注釈である。宣言が生成する nominal の値 introduction はその生成 body 内の規則である。
source の constructor application は通常の App として扱う。constructor の任意の効果付き引数式は
通常の式判断では許すが、その評価を \(\forall\) 導入の内側へ移さない。

constructor application 自体をさらに一般化できる全 source 値制限の対応は、本稿の
rank-2 payload の導出だけで閉じたとはしない。その元の受理要件を新しい拒否規則へ変更しない。

\(Q\) は通常の宣言的型付けと同じ qualified 前提である。
必要な型 instance、直接境界の有限証拠、選択証拠、role 証拠を規則木に記す。
これはユーザーに新しい annotation を要求する規則ではなく、
型推論器が将来見つけるべき数学的な証人である。
型 instance の存在と、推論器がそれを発見できることを同じ定理にしない。
一方、実行 metadata の対応表まで任意の証人として渡すことは、本稿の橋の解決とは数えない。

型境界は選択稿の直接証拠規則だけを使う。
特に optional record の適合の非推移性を保ち、二つの適合から一段の適合を捏造しない。
impl の全 body と payload 宣言は rigid binder の下で同時に検査する。
値 SCC は guarded な有限参照表、computed root は source 順の計算である。
未初期化の computed binder を読む規則はない。
この規則は linker の全経路が適切な root 環境を作るという命題を前提へ隠さない。

\section{rank-2 を source の導入と使用へ戻す}
\label{sec:poly}

\subsection{全称導入、除去と宣言形成}

\(\mathcal E\) は \(\Gamma\)、共有値 graph の型、captured scope/view の型付き環境、
自由な qualified 証拠の型をまとめた環境とする。
\[
 \frac{\Delta,a;\Gamma;Q\vdash_v n:B\mid d
       \quad a\notin\fv(\mathcal E,\D)}
      {\Delta;\Gamma;Q\vdash_v n:\forall a.B\mid\mathsf{AllI}(a,d)}.
 \tag{S-AllI}
\]
\(a\) は fresh rigid である。期待される \(\forall a.B\) が source annotation にあれば
その binder を使い、一般化の場合も同じ環境自由変数条件を使う。
追加の qualified 義務を閉じずに捨てる規則はない。
全称導入は値を構成する判断にだけ存在する。
\[
 \frac{\Delta;\Gamma;Q\vdash_v n:\forall a.B\mid d\quad\Delta\vdash T\ \mathsf{type}}
      {\Delta;\Gamma;Q\vdash_v n:B[T/a]\mid\mathsf{AllE}(d,T)}.
 \tag{S-AllE}
\]
型適用の記号は証明木の \(\mathsf{AllE}\) にあり、ソースを書き換えない。
多相値を計算して得る \(e\) に対しては、先にその計算を評価し、結果 binder に AllE を適用する。
一般の \(e\) を値に擬装する規則ではない。

宣言形成では、constructor \(c\) ごとに
\[
 \Delta,\bar a\vdash S_c(\bar a)\ \mathsf{type},\qquad
 c:\forall\bar a.\ S_c(\bar a)\longrightarrow D\langle\bar a\rangle
 \tag{S-Decl}
\]
を検査する。\(S_c\) の record/tuple/function 引数に \(\forall\) があってよい。
型参照は宣言表に登録した identity を使い、body の検査後に有限表を閉じる。
opaque な新しい型等式や無進展の alias cycle は認めない。

\[
 \frac{\Gamma\vdash_s e:S_c(\bar T)\mid d}
      {\Gamma\vdash_s c(e):D\langle\bar T\rangle\mid\mathsf{Con}(c,\bar T,d)}
 \quad
 \frac{\Gamma,x_c:S_c(\bar T)\vdash_s e_c:B\mid d_c\quad\text{各 }c}
      {\Gamma,z:D\langle\bar T\rangle\vdash_s
           \mathsf{case}\ z\ \{c(x_c)\mapsto e_c\}:B\mid\mathsf{Cases}(\bar d_c)}.
\]
多相 payload の \(\forall\) は pattern binder の環境へそのまま入る。
単相な \(S_c\) を、取り出したという理由だけで一般化しない。
record field の取得も宣言済みの field 型をそのまま返す。

\subsection{構成関数}

核への関数 \(\El_v,\El\) は規則木への再帰で定義する。
値判断を式として使う規則は \(\El(\mathsf{Val}(d_v))=\Ret\El_v(d_v)\) である。
\[
\begin{array}{ll}
 \El_v(\mathsf{AllI}(a,d))=\Lambda a.\El_v(d),&
 \El_v(\mathsf{AllE}(d,T))=\El_v(d)[T],\\
 \El(\mathsf{Con}(c,\bar T,d))=
    \Call(v_c[\bar T],\Delay\El(d)),&
\end{array}
\]
ここで S-Decl から生成する値は
\(v_c=\Lambda\bar a.\lambda^{\mathsf V,\pi_c}z.\Ret c_{\bar a}(z)\)。
\(\pi_c\) は宣言した constructor signature の source profile であり、型だけから権限を追加しない。
この値の型は payload の一変数環境で constructor introduction を適用して構成する。
source application は同じ一つの receiver を確保してから引数を force するので、
データ構築であることを理由に receipt を省略しない。
複数 payload は宣言された arity/tuple の形を使い、各実際の App を順に構成する。
生成 body の値 introduction と、ユーザーが書く constructor application を同一視しない。
Cases は scrutinee を一回評価し、constructor switch と source pattern の順序付き計画へ移す。
全称型の中身を tag 検査で開いたり、全称値を使用前に一つの単相 instance へ固定したりしない。
列挙しない constructor に対しては \(\mathsf{MatchFailure}\) を生成できる。
上の全 constructor を書いた規則表示は、その全域な特例である。

\begin{lemma}[source 全称導出の elaboration]
\label{lem:source-poly}
有限の S-Decl と上記 source 導出から、\(\El_v(d_v):A\)、\(\El(d):A\) の
核の実行型導出を構成できる。scope/view を含む自由な環境は同じものを捕捉する。
この構成は、完成した核 code が型付きであるという premise を使わない。
\end{lemma}
\begin{proof}
値/式/pattern の規則木に同時帰納する。
AllI では帰納仮説が rigid \(a\) の下の値の導出を与え、
環境自由変数条件が核の型抽象規則の前提そのものである。
AllE では核の全称除去を使い、型・row・family・qualified 義務・自由な scope template に
一つの同時代入を行う。型代入は動的 receiver を確保しない。
Con では同じ \(S_c(\bar a)\) を持つ formal から constructor introduction により \(v_c\) の型を作る。
帰納仮説の計算を delay して、その一つの instance の Call の前提にする。
Cases の各 payload binder はこの同じ template であるため、
各 arm の帰納仮説を switch の前提にできる。field の場合も同じ環境型を保持する。
通常構文の各場合は第\ref{sec:all-source}節の表の型付き穴を使用する。
宣言参照は有限表の signature を使い、body を無限に展開しない。
全 body の有限検査後に同じ表を閉じるので、未検査 impl や constructor を循環した前提にしない。
\end{proof}

\begin{lemma}[型代入と source 使用の共有]
上記導出へ capture-avoiding な \([T/a]\) を施してから elaboration することと、
elaboration 後に同じ全 packet 代入を施すことは alpha 同値である。
項の binder を既存の同じ多相値で置換する場合も、各 use の AllE は保存される。
\end{lemma}
\begin{proof}
型代入は source binder identity、constructor identity、注釈の有無を変更しない。
AllI の内側 binder は先に fresh にする。Con/Cases は同じ宣言 template への代入であり、
AllE では型代入の合成則を使う。qualified 境界では一つの同時代入を全 child に施す。
項置換の variable case は、その use に付いた AllE 列を、置換した値の外に同順で置く。
各使用時の型が異なっても、捕捉している共有値・receiver は同じ値のものを参照する。
他の構文は子への帰納仮説を結合する。以上で有限規則木の全場合を尽くす。
\end{proof}

\paragraph{ソース上の判別導出。}
数学的な宣言表示として
\(\D(C)= (\forall a.\ a\to a)\to P\) を取る。
これは新しい declaration spelling の提案ではない。
名前解決後の source
\[
 \mathsf{case}\ C(\lambda x.x)\ \{C(h)\mapsto (h\ 0,\ h\ \mathsf{true})\}
\]
について、constructor の期待 payload で \(a\) を rigid にして \(\lambda x.x:a\to a\)、
AllI により \(\forall a.a\to a\) を得る。
pattern 環境の \(h\) に同じ全称型を入れ、左の use では AllE(int)、右では AllE(bool) を使う。
従って結果は \((\mathsf{int},\mathsf{bool})\)。これらは全て source 規則の instance であり、
同じ \(h\) から二つの核の型適用を構成できる。
\(C\) が \(D\langle\mathsf{int}\rangle\) の単相 field を持つ場合には第二の AllE の前提がない。

\paragraph{reference と効果。}
\(\mathsf{Ref}(\forall a.B)\) は固定された一つの値型を格納できる。
\(\forall a.\mathsf{Ref}(B)\) を型 \(\mathsf{Ref}(a)\) の共有セルから作ろうとすると、
\(a\in\fv(\mathcal E)\) となり AllI は成立しない。
セル生成の計算を \(\Lambda\) の外へ動かす規則もない。
効果のないことだけから一般の計算を AllI に入れる規則を追加していない。
これは元の値制限の source 版であり、rank-2 のために新しい動的制限を導入したものではない。

\section{全構文の実行型を source から構成する}
\label{sec:all-source}

以下の表は \(\El\) の通常 case と、その source 前提である。
境界/選択の子には選択稿の有限規則木を使い、任意の型付き adapter を直接証拠へ昇格させない。
型注釈と全称証拠を消しても、元の source の呼出し・遅延・計算順序を保持する。

\begin{longtable}{>{\raggedright\arraybackslash}p{.19\textwidth}>{\raggedright\arraybackslash}p{.73\textwidth}}
\toprule source & 前提と生成規則 \\
\midrule\endhead
Lit & 宣言された literal 型。\(\Ret\) で返す。\\
PrimitiveOp & 固定 primitive signature を一度 instance 化し、所定 arity の callable を作る。\\
Var & 環境の型と有限 AllE/qualified-use 列。保存済み computed root を再実行しない。\\
App & callee の Function 導出、引数の導出、引数/結果の直接境界。callee を一回評価し、引数全体を \(\Delay\) にして完全 Call へ渡す。\\
RefSet & 同じ invariant payload 型の reference/update signature。reference と RHS を source 順に一回評価し、通常の protocol を呼ぶ。\\
Lambda & formal の型と source entry、parameter pattern の成功環境、body の結果導出。Value は入口で force してから pattern、Computation は保持した carrier に対する明示された pattern/force 計画を使う。source 注釈の有無は保持する。\\
Tuple & 各 component 導出を source 順に bind して tuple。\\
Record & Field/Spread の有限列。全 RHS を順に評価し、可視写像へ右優先に追加。optional absence は上書きしない。\\
PolyVariant & 同じ tag/arity と各 payload 導出。source 順の bind 後に構造導入。\\
Select & 選択稿の M/R/F 規則木と全 qualified child。\(\mathsf{BuildSelect}\) で code を構成し、receiver を通常の App へ渡す。\\
Case & scrutinee 型、各 pattern の成功環境、bool guard、結果への各直接境界。順序付き試行を生成。\\
Catch & 同じ operation instance の payload/response、raw continuation の型と各 arm の導出。shallow の通常/転送/試行計画を生成。\\
Block & source 順の環境拡張。Let は guarded 値 graph または一度の計算、Expr は結果破棄、Module は解決済み表と順序付き root。\\
\bottomrule
\end{longtable}

pattern は、Wild、Lit、Var、Tuple、List、Record、PolyVariant、Con、Ref、As、Or の11形である。
Wild は空環境、Var は入力型の binder、Lit/Ref は宣言された比較結果 bool、
Tuple/List/Record/PolyVariant/Con は実際の payload 型で子を検査する。
Record default は先行 binder 環境で検査し、欠落時だけ実行する。
As は子の成功時に元の入力 binder を追加する。
Or は両枝の同じ source name の共通部分だけに merge binder を作り、
成功した枝から直接境界を経由してその binder を埋める。
pattern に多相型があっても、これらはその値型全体を運ぶ規則である。
全ての request は中断として伝播し、失敗試行の効果を巻き戻さない。

\begin{theorem}[名前解決後の source 導出から実行型への構成]
\label{thm:source-value}
上記の有限な source 宣言・式・pattern 導出と、選択稿の有限な境界/role/選択導出から、
全称 payload を含む核の実行型導出を構成できる。
この算法は与えられた導出について停止し、未解決の義務を実行器に送らない。
\end{theorem}
\begin{proof}
有限な宣言 signature 表を先に登録する。
式/pattern の子と direct evidence の子に同時構造帰納、role の有限 DAG に高さ帰納を使う。
表の Lit/Var/PrimitiveOp は値/instance 規則、App は let、Delay、Call、明示 adapter の規則、
Lambda は formal と全 capture を同じ環境に入れた値導入である。
parameter pattern の成功環境を body へ渡し、失敗は MatchFailure、default/ref 比較の request は
そのまま伝播する。pattern の force は receipt より後の実際の entry 計画に残す。
RefSet は既出 protocol の導出である。
構造値は子の導出の順序付き bind と introduction を組み合わせる。
Case/Catch は各成功環境で得た子の導出を attempt の型付き穴へ入れる。
Or の環境は両枝で埋まる共通 binder だけなので、他方専用の未束縛 lookup は生成されない。
Block の有限 prefix への帰納で環境を拡張する。
Select は選択稿の M/R/F の局所構成、全称 case は補題\ref{lem:source-poly}を用いる。
これで13式と11 pattern の全場合が揃う。
全 body を同じ署名表の下で検査してから閉じるので、再帰 code の展開停止は要求しない。
構成算法は有限子を一回辿り、参照は既登録 label を出力するため停止する。
\end{proof}

この定理の結論は \emph{実行型}である。
source から正しい保護・解除・効果 inclusion も得たという結論を、同じ表から省略して推測しない。
そのために残る生成箇所を第\ref{sec:origin}節以降で扱う。

\section{再開する実行 owner と、過去の receiver を分ける}
\label{sec:owners}

\subsection{旧橋を否定する source 判別導出}

根拠は charter 第13節の receiver expiry と、
\nolinkurl{notes/design/2026-10-02-typed-source-owner-realization.md} 第2--3節である。
後者の具体機械は Draft だが、元の receiver を終了後に再利用しない条件は承認済みである。
以下の判別は、検索で unwind した owner を後者の機械どおり失効させる場合について述べる。
保存された control の到達可能性だけから元の receiver の生存を導く規則は、その機械にはない。
これにより既存の source-owner 機械への橋を否定するのであって、
全ての可能な suspension 意味論が同一だと仮定した反例ではない。
修正版は既存の具体化に従う候補であり、Draft の表現を新たな Authoritative 仕様へ昇格させない。
実行 owner の再確保を、過去の証拠の alpha 変換と同一視しない。

型付きの operations \(P:1\to1\)、\(Q:1\to\mathsf{int}\) を取る。
named header \texttt{my worker cb = ...} の一つの unannotated formal \(cb\) と、
その NamedHeader provenance を保持した名前解決済み AST を
\[
 \lambda cb.\ \mathsf{Block}\bigl(
      \mathsf{Expr}(P()),\quad
      \mathsf{Catch}(cb(),\{Q(),k\mapsto0;\ v\mapsto v\})\bigr)
 \tag{W}
\]
とする。\(cb\) の invocation は \(Q\) を出すものとし、
外側の handler \(H\) が \(P\) を処理して、受け取った raw continuation を再開する。
全て元稿の13式で書ける有限 source である。
単独の Anonymous lambda に書き換えた場合にも同じ source frame が選ばれるとは仮定しない。
この一引数 named group では cb の introduced/terminal frame は同じであり、
第\ref{sec:named}節の direct-call guard と同じ手順で cb の scope を生成できる。
ここでは実機での実行結果を主張しない。

通常呼出しの owner を \(u\)、callback 受領での保護 Boundary を \(b\) とする。
omission の \(\mathsf{Admit}_b=\varnothing\)、\(b.\mathsf{receiver}=u\)。
\(P\) は callback の effect port の外で発生するため、当該 \(b\) はその発生を保護しない。
\(P\) を外側の \(H\) が選ぶ時点で \(u\) の scope は request と一緒に切り出され、
現在検索中の context から除かれている。\(u\) はその後自動的に active へ戻らない。
再開して残りの Catch を実行する owner は fresh \(u'\) である。

\begin{proposition}[receiver の一様改名は source 寿命を保存しない]
\label{prop:old-failure}
旧機械の規則に従って saved packet 内の \(b.\mathsf{receiver}=u\) を \(u'\) へ改名すると、
(W) の再開後の \(Q\) の選択は、元の receiver expiry 規則と異なる。
従ってその改名を用いる旧 source-to-machine 対応は成立しない。
\end{proposition}
\begin{proof}
再開した Catch の handler \(h\) は owner \(u'\) を持つ。
callback の typed path とその実行時 Observe は実際に存在する。
旧改名では \(b.\mathsf{receiver}=u'\) で active となり、
\(P(Q,h)=1\)、\(G(Q,h)=0\)。よって \(h\) は \(Q\) を処理できない。
source-owner 機械の unwind/enter 規則で元の receiver を固定すると
\(b.\mathsf{receiver}=u\)、\(\neg\Act(u)\) であるため、
この唯一の保護 witness は incidence に入らない。
同じ operation/instance と active owner の \(h\) は通常規則で \(Q\) を処理して0を返す。
raw continuation の現在 owner への Receive は \(b\) 自体を作り直さない。
差は同じ family の判定や値型からではなく、過去の receiver を変更したことだけから生じる。
両者の計算は実行型 \(\mathsf{int}\) を持ちうるので、値型保存だけではこの差は発見できない。
\end{proof}

\subsection{修正版の具体的な状態と遷移}

\(\Core_{\rm life}\) では値・型・通常 code・ordinary receipt・View/Flow/Observe は旧稿と同じ。
次だけを置き換える。過去の evidence identity と実行 control slot を別 sort にする。
Boundary receiver、Receive owner、request origin、既存 view、scope \(s\)、target \(t\) は
実際に確保された immutable な値であり、後の resume の束縛変数ではない。
source code に残る未実行の allocation binder は通常どおり alpha 変換できる。
この区別は静的 label と動的 identity の区別であって、型を二重化するものではない。

saved context は7 frame の列に、各 scope の元の owner と source descriptor を保持する。
再開時の control 環境を \(\rho\) とし、outer frame から inner frame へ次を実行する。
\[
\begin{array}{ll}
 \mathsf{enter}(i,u,C)=(u,\mathsf{borrow},C),&\Act(u,C),\\
 \mathsf{enter}(i,u,C)=(u',\mathsf{own},C+u'),&\neg\Act(u,C),\quad u'\text{ fresh}.
\end{array}
 \tag{L-Enter}
\]
該当 control slot を得られた owner に結び、以前の slot 値と completion mode を frame に保存する。
borrow の完了は active 集合を変えず、own の完了はそこで確保した \(u'\) だけを除く。
child が返ったら親の suffix を実行し、raw continuation の root だけが resumer へ返る。
選ばれた shallow handler は再設置しない。
転送により保存された handler は、その source descriptor から新しい activation を作り、
\(\rho\) により解決した実行 owner を使う。
既存の live handler と混同したり、body の外まで capture 範囲を拡大したりしない。

\textbf{過去の evidence への作用は恒等。}
L-Enter は過去の \(b.\mathsf{receiver}\)、Receive、Observe、view の値を変更しない。
保存した binding を現在の body が受領する時は通常の
\(\mathsf{Receive}(\rho(i),\mathsf{slot},v)\) を追加するだけで、Boundary を追加しない。
新しい通常の source Call は、その source receipt 規則を実際に再実行するので fresh Boundary を作れる。
単なる raw resume は call entry を飛び越えた suffix から始まるため、その receipt を再実行しない。
共有資源の read は、stored evidence identity を固定した値を現在の typed path へ transport する。

\textbf{release 状態。}
既に確保した \(s,t\) と \(L\subseteq Scope\times Target\) も過去の値として固定する。
各 resume が読むのは現在の \(L\) であり、creation-time snapshot を上書きしない。
\(L\) を新しい control owner 用に複写する規則を廃止する。
未実行の source target introduction に新たに到達した場合だけ fresh \(t'\) を確保する。
RC と RP は元の receiver の activity を見る同じ式である。
その receiver が dead なら、古い release entry と古い protection の双方が有効性を失う。
保存した target の同一性と、source の新しい通常呼出しで作る別 target の同一性を分ける。

\begin{lemma}[寿命、借用、複数回の再開]
\label{lem:life}
\(\Core_{\rm life}\) の各有限 prefix について、
(a) 過去の receiver が L-Enter だけで復活しない、
(b) borrowed completion が外側の live owner を終了させない、
(c) 二回の resume が saved code/evidence を変更しない、
(d) 同じ live receiver/target は現在の \(L\) を共有する、が成り立つ。
\end{lemma}
\begin{proof}
初期状態に既存の違反はない。一段の最後の規則で場合分けする。
L-Enter の active 分岐は既に active な同じ \(u\) を借りるだけである。
inactive 分岐は全旧 support と不等な \(u'\) を追加するので、旧 \(u\) の activity を変えない。
完了は保存した mode を検査し、own の場合だけ自分の新しい点を除く。
ordinary Call/receipt は実際の source 呼出しに対応し、旧 identity を再利用しない。
metadata transport は旧点を保持する Flow、release は RC の同じ \((s,t)\) の追加だけである。
resume は frame と control 環境を新しく作り、saved code/evidence に書き込まない。
共有 read は現在の値を返し、過去の状態への rollback 規則はない。
各場合で(a)--(d)が保たれるので、prefix 長の帰納法で結論を得る。
\end{proof}

\subsection{保存・進行と有限な効果表示への戻し方}

\begin{theorem}[修正版の実行型保存と進行]
\label{thm:life-value}
核の実行型導出を持つ閉じた計画は、\(\Core_{\rm life}\) でも各有限実行 prefix で型を保存する。
値、型に合う request、列挙した failure、次の遷移のいずれかであり、\(\Bad\) に到達しない。
\end{theorem}
\begin{proof}
旧稿の値型保存は任意の全域 eligibility に対して示されている。
型と family/resource identity を変えない L-Enter/leave は同じ code の型導出を使う。
active 分岐と inactive 分岐は排他的かつ全域であり、有限 support の外に fresh 名がある。
保存 frame の各穴の型は変わらず、親 suffix を let の規則で合成する。
handler の owner slot と activation は型の payload/response instance ではないため、
forwarded handler の再導入も同じ構文的 arm 型を持つ。
raw continuation の response には同じ項置換補題を適用できる。
過去の view を改名しないことは underlying constructor を変えない。
release の状態参照も値型を変えず、RC/RP の query は有限 DAG 上で全域のままである。
その他の redex は旧稿と同じ。従って一段保存と canonical forms による進行が得られ、
有限 prefix 長への帰納で結論する。
\end{proof}

精密な効果定理には、旧稿で検査済みの小さい row をそのまま使えない。
修正版では旧稿の有限 complete cell に、各 register の
\(\mathsf{evidence},\mathsf{control},\mathsf{viewOccurrence}\) という action sort と
frame completion mode を加える。追加は有限個の列である。
L-Enter の transfer は次の有限 relation として作る。
\begin{enumerate}
\item evidence register と既存の \(s,t,L\) は恒等。control が指す旧 \(u\) の activity を読む。
\item active なら同じ点を control 出力へコピーし、mode を borrow にする。
inactive なら旧 support と不等の fresh 点を出力し mode を own にする。
\item 現在の Receive と forwarded handler は出力 control 点を使う新しい行。
旧 Boundary の receiver 列と Admit はコピーするだけである。
\item 新しい executing View occurrence は fresh とし、typed packet と effect port は保持する。
新しい emission は現在実行中の全 View を観測し、過去の Observe は固定する。
\item leave は mode に従い、own の点だけ activity を落とす。
hidden な点の一致、mode、route が保持されない場合は整合する全候補へ戻す。
\end{enumerate}
全 input tuple は共同 complete cell であり、引数/結果/ambient の相関を独立に選ばない。
旧稿の Boolean region 像、fresh、forget を使うので、各 relation は有限に列挙できる。
未知の incidence を都合のよい許可に確定する規則はない。

\begin{lemma}[修正した transfer の表示保存]
\label{lem:life-transfer}
具体状態を表示する入力 cell から \(\Core_{\rm life}\) の metadata step を一つ行うと、
得られた状態の cell は上記 relation の出力集合に含まれる。
\end{lemma}
\begin{proof}
具体的な activity が真なら列挙された active 分岐、偽なら inactive 分岐が選ばれる。
前者では出力の control identity が同じ点、後者では全旧点と不等であり、fresh の式を満たす。
evidence の各座標は具体的にも変わらない。
新しい Receive/handler/View は記述した source 操作の式と同じ点を使う。
leave は保存した具体 mode に合う一方の分岐にある。
hidden の一致・route が失われた場合も、その具体的な refinement は全候補の一つである。
RC/RP は既存 \(s,t\) と現在の activity に対する同じ Boolean 式なので、旧稿の座標ごとの証明を使える。
ただし旧稿の bound-release-state 複写の場合は削除し、恒等 read の場合へ置き換える。
これが変更した全操作であり、残りは旧稿の表示保存を使う。
\end{proof}

\begin{theorem}[修正版に対する精密な型健全性]
\label{thm:life-full}
\(\Core_{\rm life}\) の生成規則と上記 transfer を用いて、全 mandatory edge、
cut-context summary、全共同入力の inclusion を満たす有限な \(A!\Phi\) 導出があるとする。
その計画の全有限実行 prefix で値型と immediate/latent/reentry 契約が保存される。
原 receiver の失効は補題\ref{lem:life}に従い、空の outward 契約から request は出ない。
\end{theorem}
\begin{proof}
実行型は定理\ref{thm:life-value}。
cut の7 frame のうち owner entry/leave と saved view の操作を上の relation へ置き換える。
各 frame の残りの let/handler/mapop/wait/release の型付き穴は同じである。
frame 列の長さへの帰納により、具体 suffix は修正した cut summary に表示される。
finite relation image は入力集合について単調なので、有限 powerset の固定点と置換補題を使える。
一段の具体 metadata は補題\ref{lem:life-transfer}、query は完全 cell 上の同じ Boolean 式である。
従って新しい owner の下で旧 protection が無効になって handler が選べる場合も、
修正した graph の handler 分岐に含まれる。旧 filter の結論を流用しない。
No の試行から戻る場合は current state を再観測し、wait の pending query をその時の cell で検査する。
全 mandatory edge の inclusion から次の入力 port の表示を得る。
これで精密な状態判断の一段保存が得られるので、prefix 長に帰納する。
\end{proof}

\textbf{依存の意味。}
この定理は旧稿の値型・有限 cell・query・cut の紙上補題を依存に持つ。
独立査読の代用ではない。機械を修正したため、旧稿の具体的な効果導出と
選択稿の \(\mathsf{incl}\) 証拠は新しい relation で再検査する。
旧機械で成り立った小さい residual が同じであるとは主張しない。

\section{注釈から作る origin inventory と target}
\label{sec:origin}

\subsection{型の同一性と、注釈の由来を別に保持する}

型注釈の lexical row binder に identity \(\epsilon\)、その occurrence に \(j\)、
boundary を導入する source node に \(b_0\) を付ける。
静的 origin key は \((\epsilon,j,b_0)\) である。
\(\epsilon\) と \(\epsilon'\) に同じ型を代入しても、二つの origin key は等しくならない。
型を \(\forall\) の内部へ運ぶ際も、その binder の scope と shadowing を先に解決する。
\texttt{'e?} の lexical name は、同じ scope で解決された \(\epsilon\) を参照する。
単に末尾の文字を落とした row ID で動的 receiver を探索する操作ではない。

source 導出には型・効果 port の有限 template graph を作れる。
各 node は source occurrence と型の child path、辺は通常の alias、構造の prefix/除去、
同じ直接境界の correspondence、Function の正負の port、scheme の一つの instance map である。
再帰は同じ template label への参照とし、実行時の Flow DAG と区別する。
型の shape が遅れて確定する位置には、その元の occurrence に属する template を残す。
無限の nominal 展開や、同じ型だから同じ由来という逆向きの辺を作らない。

\subsection{annotation のレシピを生成する規則}

\begin{longtable}{>{\raggedright\arraybackslash}p{.29\textwidth}>{\raggedright\arraybackslash}p{.63\textwidth}}
\toprule source occurrence & 生成するレシピ \\
\midrule\endhead
反変 concrete \([C,X_1,\ldots,X_n]\) & 実際の入力 port と各 symbolic binder を保持し、concrete \(C\) の Admit と profile を作る。actual inflow へ \(C\) の seed は追加しない。\\
共変 concrete row & 同じ結果 port の inclusion。新しい保護や捕捉許可を作らない。\\
callback wildcard & 既存の view を運び、新しい Boundary を作らない。\\
明示 computation formal & carrier/force の source entry を保存する。wildcard は当該 direct port の AllowAll、具体項はその局所 Admit。\\
注釈のない formal & formal の source binder に seed origin を保持する。下記の実際の upper-use exposure だけで新しい profile を作る。\\
既知の外部 Name & 元の値と既存 origin の transport。名前使用が無注釈であることを新しい formal seed としない。\\
全称導入/除去 & 静的 recipe binder を capture-avoiding に運び、全 packet に同じ instance を使う。既に存在する動的 receiver を作り直さない。\\
\bottomrule
\end{longtable}

polarity は Function の引数辺で反転し、結果辺で保持する。
カンマ、セミコロン、許された改行は同じ順序付き row items に正規化する。
concrete atoms を集めるときに型変数を数えないことと、その binder/依存辺を消すことは異なる。
normalization は source span の対応を保ち、別の \(\epsilon\) を統合しない。

\textbf{省略位置の source 条件。}
2026年10月6日の方向付き決定を、解いた型の全走査へ置き換えてはならない。
既に保護 seed を持つ formal/view \(v\) が、実際の source App により
\(v<:F(A,E^-,E^+,B)\) の upper use \(u\) に露出した場合、
その \(u\) の \(E^+\) にだけ seed の対応を付ける。
provider の lower bound から逆向きに \(E^+\) の保護を配らない。
元の finite source-use graph でこの生成を反復する場合は、
\((\text{元 seed},\text{source use},\text{template port})\) の組ごとに一度登録する。
recursion のためだけに source occurrence を無限に複製しない。

この規則の\emph{局所構成}は、source use と seed の有限な組に対して全域である。
しかし、任意の source inference の replay/一般化が元 seed を保持することまで、この列挙からは導けない。
source 導出に既存の alias/structural edge を明記することは通常の型付けであるが、
\(\mathsf{ProtectedVarAt}\) を理由なしの自由な仮定として加えることは、全言語対応の解決にはならない。
その exact producer を次の inventory の入口として残す。

\subsection{実行時の inventory は receipt から作る}

各 recipe に対応する actual source boundary-entry が実行されたとき、
通常の \(\mathsf{Boundary}(b,r,v,\pi)\) を作る。
その profile に属する row component occurrence \((\epsilon,j)\) について
\[
 s=\mathsf{scopeKey}(b,\epsilon,j),\qquad
 \mathsf{Origin}(s)=(b,r,\epsilon,j)
 \tag{O-Intro}
\]
を保持する。\(\mathsf{scopeKey}\) は sort を分けた単射でよく、
同じ Boundary の同じ occurrence を二回引いても同じ key を返す。
regular/nominal template が実際の path を要求する場合、\(j\) は元の node とその有限 path を含む。
全潜在 path を receipt 時に列挙する必要はなく、実際に参照された有限 path の key を使う。
これ自体は新しい Boundary や grant を作らない。

inventory は各 typed view に沿う route の記録である。
\[
 I(v)=\{(s,\epsilon,j,\tau)\mid
       \tau\text{ は O-Intro から }v\text{ の該当 port への実際の typed route}\}.
\]
実装上の集合の全列挙を要求せず、有限 DAG と各 edge の prefix template で保持できる。
各実行 prefix で query が辿る DAG と observed path は有限である。
生成規則は O-Intro、実際に実行した Flow の合成、source が許した projection に限る。
closure construction は latent view に inventory を捕捉し、実際の call/force 前に Observe しない。
同じ pointer/family/解いた row 型だけから inventory entry を追加しない。
shared reference の read、case binder、構造 adapter、dictionary field の取得も、
その操作の実際の typed path を使う。

\begin{lemma}[inventory の構成、健全性と経路についての完全性]
\label{lem:inventory}
O-Intro と上記 transport が作る任意の有限 prefix で、
各 inventory entry は元の一つの実際の boundary-entry と typed route を持つ。
逆に、その prefix のこれらの規則だけで構成した route は inventory の query に現れる。
元の receiver は transport、AllE、resume によって変更されない。
\end{lemma}
\begin{proof}
prefix 長に帰納する。O-Intro はその transition の \(b,r,\epsilon,j\) を唯一の起点にする。
Flow の追加は、新しい endpoint へ到達する既存の route にその一辺を連結する。
projection の場合は対象 component の prefix に一致する route だけを残す。
逆に新しい endpoint へ到達する route の最後の edge は、この追加された一辺であるため、
前半は帰納仮説により既存 inventory にある。branch は実際の成功枝の endpoint を使う。
AllE は型/静的 binder の一様な解釈だけを変え、既存の \(b,r,s\) を固定する。
resume は第\ref{sec:owners}節により過去の evidence を固定する。
この生成規則を使わない新しい値は空 inventory から始まるので、不正な起点を作らない。
DAG の逆向き走査と有限 path の検査は停止する。
\end{proof}

これは\emph{実行された source 境界と transport}についての完全性である。
全 source から正しい初期 recipe が生成されるという主張に、量化対象をすり替えない。
特に型の解からだけ protected seed を再構成する方式には、この補題を適用できない。

\subsection{target の生成と、未確定な複数 occurrence の意味}

印 \(m\) の row binder を \(\epsilon_m\)、実際の output occurrence を \(j_m\) とする。
印の対象入力 occurrence を選ぶ静的関係を \(A_m\) と書く。
これは source の解釈が決める有限の occurrence/template の集合であり、
動的 receiver 名、効果の family 名、解除後の handler 選択結果を含まない。
recipe は実際の output slot に fresh target \(t_m\) と delimiter を作り、
到達した inventory のうち
\[
 (s,\epsilon_m,j,\tau)\in I(v),\qquad j\in A_m
 \quad\Longrightarrow\quad
 \mathsf{Target}(t_m,s,\mathsf{Origin}(s).r,j_m)
 \tag{O-Target}
\]
を target view に結ぶ。
表は inventory の有限 query で構成され、完成した scope/receiver 表を型付けの前提として渡さない。
必要な correspondence は保持した template であり、潜在 event を先に実行するものではない。
後から同じ target に達した実際の route は同じ query で読まれる。

RC が更新するのは \(A_m\) と From の両方を満たし、かつ元 receiver が active な pair だけ。
出力 whole packet に RP を適用するため、\texttt{'f} 内の同じ origin も同じ \((s,t_m)\) を読む。
元の input view、別 target、別 lexical binder の保護は変更しない。
同じ marker の単なる再観測では target を作り直さない。
source の別の通常呼出しで同じ output occurrence に新しく入る場合は fresh target である。

\begin{lemma}[固定した occurrence 選択からの target 構成]
\label{lem:target}
source の \(A_m\) が固定され、recipe の起点が実際の boundary-entry から生成されたなら、
O-Intro/transport/O-Target は旧稿の well-scoped origin/target 条件を構成する。
\(\Core_{\rm life}\) では、その target の release は元の receiver の寿命を保つ。
\end{lemma}
\begin{proof}
補題\ref{lem:inventory}が各 \(s\) の唯一の起点と元の \(r\) を与える。
O-Target は \(r\) をその起点から読み、印と同じ lexical binder および選択された occurrence を検査する。
従って未束縛 scope、別の receiver、family 一致だけによる target は作られない。
型/証拠の代入は全 occurrence と template へ同じ写像を適用し、既存の動的値は固定する。
新しい output entry だけが fresh target を作り、transport は同じ target を保持する。
RC の From は現在の有限 route によるので、解除後の eligibility に依存しない。
最後に補題\ref{lem:life}を用いれば、終了済み receiver の解除権限を復活させない。
\end{proof}

\textbf{ここで残る source の意味。}
同じ \(\epsilon\) を持つ独立な ingress が二つある場合、
\texttt{'e?} の字面だけから「両方を対象にする」か「一方の occurrence を選ぶ」かは、
単一 ingress の規則では決まらない。
例えば数学的な scheme
\[
 (()\to[\texttt{'e}]A)\to(()\to[\texttt{'e}]B)
       \to[\texttt{'e?}](A,B)
\]
について、型変数の共有から二つの Boundary identity の同一性を導くことはできない。
一方、別の \(\epsilon'\) が偶然同じ row 型になったことを理由にその保護まで解除する案は、
既存の独立 scope 保存に反する。
複数入力の \(A_m\) の選択はユーザーへ照会中であり、回答前に一般の source 規則として採用しない。
補題\ref{lem:target}はこの一箇所を明示した構成であり、全 source 対応の完了とは呼ばない。

\section{compose を実際の呼出し構造で追う}
\label{sec:compose}

承認済みの期待は次である。
\begin{Verbatim}[fontsize=\small]
my compose f g x = f (g x)
compose : ('a ['e] -> ['f] 'b)
       -> ('c -> ['e] 'a)
       -> 'c -> ['e?, 'f] 'b
\end{Verbatim}
この表示の主型性と一般推論はここでは使わない。
まず実行構造について導出できることを、origin の導入から分ける。
元の IR は \(\mathsf{Lambda}(Pat,Expr)\) と二項 App であり、三つの formal は入れ子の lambda になる。
第\ref{sec:all-source}節の規則は、最内の body を次へ翻訳する。
\[
 \Call(f,\Delay(\Call(g,\Delay(\Ret x)))).
 \tag{C-Body}
\]
両者が値 binder であるため、callee の administrative let は省略表示した。
\(f\) が Computation entry なら、inner \(g\) は \(f\) の body がその carrier を force したときだけ実行する。
\(f\) が Value entry なら、\(f\) の receipt の後にその carrier を force する。
caller 側で \(g\,x\) を先に評価してから \(f\) を呼ぶ展開ではない。

印がこの最内 source result slot にある場合、target delimiter の位置は
\[
 \Scope_{r_x}\bigl(\mathsf{release}_{t}
     (\Call(f,\Delay(\Call(g,\Delay(\Ret x)))) )\bigr).
 \tag{C-Output}
\]
\(f\) の入口で作った receipt の receiver を \(r_f\) とすると、通常の Call により
\(\mathsf{release}_t(\Scope_{r_f}(\cdots))\) の順になる。
外の印を \(\Scope_{r_f}\) の内側に動かすことは、この source output の翻訳ではない。

\begin{proposition}[compose の output における receiver 条件]
\label{prop:compose-lifetime}
(C-Output) で \(f\) の内側から request が target delimiter へ達するとき、
\(r_f\) は既に保存された continuation 側へ移っている。
従って \(r(s)=r_f\) である scope について、target での RC の activity 条件は偽である。
\end{proposition}
\begin{proof}
request は内側から外側へ一段ずつ bubble する。
\(\Scope_{r_f}\Req(q,p,K)\) の scope 規則が、\(r_f\) を現在の context から除いて
continuation に同じ suffix を保存してから、外の release frame が request を受け取る。
release frame 自体は receiver を導入しないので \(\Act(r_f)\) は偽のままである。
内側の handler が最初の request を消費した場合には、そもそもその request について RC は実行されない。
正常 return の場合も \(r_f\) は終了し、return 自体は \(L\) を更新しない。
\end{proof}

さらに通常の currying では、\(f\) を受け取った第一の呼出し、\(g\) を受け取った第二の呼出しは、
後続 lambda を返して既に終了してから、\(x\) の呼出しが実行される。
captured view に元の Boundary を残すことは、それらの元 receiver を active にすることではない。
新しい \(r_x\) への Receive だけで古い Boundary の receiver を \(r_x\) へ変更できない。

\textbf{旧稿の局所計算の適用条件を戻す。}
旧稿は「同じ \(s_e\) が \(g\,x\) から \(f=id\) の identity edge を通り、
元 receiver が live なら、\texttt{'f} 内にも release が届く」と計算した。
identity edge の保存は正しい。しかし命題\ref{prop:compose-lifetime}により、
その \(s_e\) を単に \(f\) の入口 receipt と同一視して source 全体の証明を閉じることはできない。
\(r_f\) が dead なら保護も既存の activity filter で無効になるが、
それは「承認された \texttt{'e?} の起点と寿命を構成した」という結論ではない。

2026年10月15日の native 実験は、別の具体的な候補を既に持つ。
\(g\) の ResultEffect に \(\mathsf{ProtectBody}(E,u_g)\)、
\(f\) の ArgumentEffect に \(\mathsf{StampInput}(u_f)\)、
\(f\) の ResultEffect に \(\mathsf{ProtectBody}(F,u_f)\)、
\(f\) の Apply result に \(\mathsf{Release}(E)\) を置く。
同じ authenticated route では、identity が input tag を保つため新しい \(F\) の追加を抑え、
whole output の \(E\) を取り除くという代数計算になる。
これは局所的に査読された構成であり、捨てて最初から別の代数を作る理由はない。
ただしその \(E,F\) は native の source region/scope identity であり、
本稿の動的 Boundary receiver \(r_f\) と同じであることは示されていない。
その実験の source 全体・handler subtraction・一般寿命の証明は元の記録でも OPEN である。

\textbf{正確に残る生成義務。}
入力は C-Body の通常の source-use graph と shared assignment。
出力は (i) \(E\) の実際の保護導入箇所、(ii) 元 receiver とその生存区間、
(iii) \(u_f\) の input stamp と \(F\) の条件付き導入、
(iv) source output にある release target の対応である。
次節では直接の named-group case に対してこれらを生成する。一般の source では、生成してから補題\ref{lem:inventory}と\ref{lem:target}を適用する必要がある。
\texttt{compose} や \texttt{id} の名前を特別扱いする規則、source slot の移動、
終了した receiver の再束縛は、この義務の解決として認めない。

\section{直接の named group では compose の起点を構成できる}
\label{sec:named}

\subsection{後から発見した既存結果を、元の義務へ適用する}

次の既存結果を使う。
\nolinkurl{notes/theory/2026-10-15-callback-named-frame-source-bridge.md} の C1/C2 は
named header の実際の frame 構成を調べ、C3 は non-Var callee の一般化を反駁している。
\nolinkurl{notes/theory/2026-10-15-callback-identity-receiver-construction.md} は、
選択された frame と元の formal ごとの key、その source output での discharge を特定している。
これらは局所的に独立査読された構成であり、その premise と反例を保持して使う。

固定 Oracle の実際の所有箇所は
\nolinkurl{crates/infer/src/lowering/expr/tail.rs} の
\nolinkurl{unannotated_local_callee_return_effect}（741--798行）、
\nolinkurl{unannotated_call_frame_index}（801--828行）、
\nolinkurl{local_callee_def}（845--850行）、
\nolinkurl{lambda_predicate_subtracts}（1058行以降）と
\nolinkurl{lambda_output_predicate_vars}（1092行以降）である。
単に ``source が証拠を供給する'' とせず、次の guards を全て確認する。

\begin{lemma}[compose の静的 owner は最後の formal frame]
\label{lem:compose-static}
source \texttt{my compose f g x = f (g x)} の二つの直接呼出しは、
上記 constructor の保護 key の owner としてともに \(x\) の frame \(X\) を選ぶ。
\(g\) 由来の key は \(k(X,g)\)、\(f\) 由来の key は \(k(X,f)\) であり、相異なる。
matching output は \(X\) の body 全体の値/effect 出力である。
\end{lemma}
\begin{proof}
header の順序付き formal は \((f,g,x)\)。named group の frame range を
\([a,a+3)\) とすると、body の current Defined frame は \(a+2=X\)。
両 callee は実際に \(\mathsf{Var}\)、参照先は元の \(\mathsf{Def::Arg}\)、call-return annotation は
Unannotated である。この三つの guard は型の解から推測せず、source の構文と binder による。
\(f,g\) の introduced frame は \(a,a+1\) であり、同じ active skeleton の
\texttt{introduced < before\_frames} を満たさない。
body 内に別の named group はないので、crossing guard は偽で、選択結果は current frame \(X\)。
\(X\) は Defined なので最後の guard も通る。
実際の map key は \((X,\mathsf{DefId})\) であり、初回に確保し同じ組では再利用する。
\(f\ne g\) なので二つの key は異なる。
同じ frame の \texttt{subtracts} に対応する Pop が登録され、body 出力に適用される。
\nolinkurl{lambda_output_predicate_vars} は値と effect の両方にその列を適用するので、
\(f\) の actual result もその内側にある。
\end{proof}

この補題は、(C-Output) の live receiver を内側の \(r_f\) から得ようとした試みを修正する。
選ばれた source owner は \(X\) であり、その\emph{現在の実行}の \(r_X\) を使う。
\(f\) や \(g\) を最初に受け取った過去の curried call の receiver を復活させることではない。
静的 key と動的 receiver はまだ別物なので、次にそれを作る実際の entry 規則を与える。

\subsection{source key から実行 recipe を作る}

これは上記 constructor に対する具体化候補であり、新しい source 記法や実装上の特別扱いではない。
finite named header を走査して ordered formal IDs と terminal body label を保持する。
本文にある各 eligible direct call について、上の frame 選択を実行し、
\(\mathsf{Own}(X,p)\) という recipe を一つ登録する。
同じ source key の use はその recipe を参照する。
型の一致、family 名、関数名 \texttt{compose}/\texttt{id} を guard にしない。

terminal source lambda \(X\) の通常 entry では、現在の実行 owner \(r_X\) と
その実際の captured/formal binding から recipe の各 \(p\) の view を受領する。
\(\mathsf{Own}(X,p)\) についてだけ fresh Boundary \(b_{X,p}\) と
scope \(s_{X,p}\) を作り、その origin に \(r_X\) と元の source key を記す。
これを \(\mathsf{EnterGroup}\) と呼ぶ。新しい実行 owner は増やさない。
これは source key によって予約された一つの boundary-entry であり、
任意の captured binding に境界を追加する規則ではない。
昔の Boundary を改名する操作でもない。\(p\) の ordinary call site では、
同じ source key に属する出力 view を使い、既知の外部 callee に新しい key を割り当てない。

recipe の scope は source の\emph{当該 upper-use の出力}に付く。
共有する provider の lower graph を書き換えず、当該 consumer view からの route でだけ読む。
\(g\) の output port が source 導出で \(\epsilon_e\) に結ばれていることにより、
\(s_E=s_{X,g}\) の source attribution が得られる。
\(s_F=s_{X,f}\) は別の key であり、\(\epsilon_e\) の marker の対象にはしない。
\(X\) の出力には、この同じ inventory を引く delimiter を \(\Scope_{r_X}\) の\emph{内側}に置く。
これが実際の \(X\) body output であり、外の slot を移動したのではない。

\textbf{source 条件を満たさない場合。}
過去の receiver の失効を避けるために、任意の source scope を \(r_X\) に付け替えてはならない。
この entry を使える根拠は \(\mathsf{Own}(X,p)\) を生成した元の named header と direct use である。
一般の alias、anonymous Lambda、異なる named group からの capture には、
その source owner の構成が別に必要である。これらの source を言語から除外したことにはしない。

\subsection{入力由来を失わない二つの具体 transfer}

既存の scoped algebra の \(\mathsf{StampInput}\) と \(\mathsf{ProtectBody}\) を、
実行上の一つの寄与 \(w\) に適用する。
\(T(w)\) は実際に通った入力 route の tag、\(U(w)\) はその寄与の保護 scope である。
family/event の equality から \(T\) を追加する規則はない。
ordinary Call の carrier に有限な stamp template を付け、
その carrier が実際に force されるか、その対応した latent port が実際に実行されるときだけ、
その typed contribution に input tag \(u\) を付ける。
未実行の carrier は event も stamp された event も生成しない。

\[
\begin{array}{l}
 \mathsf{Stamp}_u(T,U)=(T\cup\{u\},U),\\
 \mathsf{Protect}_{s,u}(T,U)=
 \begin{cases}(T,U\cup\{s\})&u\notin T,\\(T,U)&u\in T.\end{cases}
\end{array}
 \tag{N-Transfer}
\]
\(u\) はその実際の source-use entry の fresh input tag であり、型適用だけでは作らない。
\(\mathsf{Protect}_{s,u}\) は profile の元の許可とは別の、omission 起源の protection 条件である。
Admit/Grant は追加せず、既にある \(s\in U\) も stamp の存在だけで削除しない。
出力の type/source key を共有していても、入力由来の寄与を新たな body 起源とみなさないための条件である。

これらは actual typed route 上の View 操作であり、全候補の Observe を消す操作ではない。
complete CallView で観測した同じ寄与にも、同じ共同入力 packet の \(T\) を使用する。
carrier 由来の route と body 由来の route を別々に選び、片方の unstamped alternative を
同じ出力へ無条件に追加する規則を置かない。
Source 型の結果が入力に由来する場合、return/field/identity はその実際の view を運ぶため \(T\) を保持する。
一般の user cast が新しく作った値へ、同じ型であることだけを理由に \(T\) をコピーしない。

この concrete 操作を持つ拡張を \(\Core_{\rm group}\) とする。
有限 cell に input-tag sort と \(u\in T\)、\(s\in U\) の共同座標を追加する。
Stamp は tag bit の OR、Protect は
\[
 [s\in U']=[s\in U]\lor\neg[u\in T]
\]
である。保護 query は、この値と元 receiver の activity、RP の同じ条件を使う。
保有しない入力相関は全 refinement へ広げる。
これらの metadata action を新しい opaque host primitive にはしない。

\begin{lemma}[group entry と入力 tag の保存]
\label{lem:group-transfer}
EnterGroup、Stamp、Protect は underlying value、operation instance、payload/response 型を変えず、
有限 relation がその具体操作を表示する。
保護と source key の対応は ordinary Call、typed transport と原 receiver expiry の下で保存される。
\end{lemma}
\begin{proof}
EnterGroup は通常 source entry の現在 owner を読み、元の \(\mathsf{Own}(X,p)\) ごとに
fresh Boundary と scope を作る。\(p\) は実際の環境 binder なので typed view が存在する。
値には新しい consumer view を作るだけで、raw constructor は固定する。
stamp は実際の force/call の typed route に一つの tag を加える。
Protect はその同じ packet の bit に上式を適用するだけである。
従って有限表示は入力の共同 cell の各座標で等式を直接計算したものになり、
hidden な相関も実際の refinement を含む。
Source key に対応する \(s,r_X\) は transport で固定される。
新しい通常の \(X\) 呼出しでは fresh \(r_X,b,s\) が導入され、過去の view を変更しない。
raw resume は過去の receipt を固定するため EnterGroup を再実行しない。
未実行の通常 Call に再開後初めて到達すれば、そこで初めて新しい entry を行う。
最後に補題\ref{lem:life}で expiry と borrowed owner の場合を処理する。
\end{proof}

\begin{theorem}[group recipe を加えた実行健全性]
\label{thm:group-safety}
上の新しい有限 relation を含めて全効果導出を生成・検査した
\(\Core_{\rm group}\) の計画は、定理\ref{thm:life-full}と同じ実行健全性を持つ。
source の直接の named-group recipe は、完成した Boundary/target 表を入力せずにこの計画を生成する。
\end{theorem}
\begin{proof}
新しい操作は underlying code の型を変えないため実行型の保存・進行は同じ。
EnterGroup の有限個の確保は actual source entry の一段の metadata 構成に含める。
Stamp/Protect は既存 View frame の有限 transfer として合成でき、cut の新しい種類は要らない。
補題\ref{lem:group-transfer}を一段表示の追加 case にし、frame 列と実行 prefix に帰納する。
Source recipe からの生成では、静的 header/use 走査で \(\mathsf{Own}\) を作り、
実際の entry で \(r_X,b,s\)、source output で target を作る。
どの rule も解いた row equality や未知の runtime table を要求しない。
全 mandatory graph edge を追加してから inclusion を検査するので、
旧機械の effect 証拠をそのまま移す循環もない。
\end{proof}

\subsection{compose id g への適用}

二つの direct use の source label を \(j_g,j_f\) とし、
命題\ref{prop:compose-lifetime}の内側 receiver \(r_f\) と、今回の \(r_X\) を区別する。
\(X\) の通常 entry は \(s_E=s_{X,g}\)、\(s_F=s_{X,f}\) を導入し、
\(j_g\) の output attribution を \(\epsilon_e\) へ、\(j_f\) を \(\epsilon_f\) へ接続する。
これはソース導出の \(g:'c\to['e]'a\)、\(f:'a['e]\to['f]'b\) の port 対応による。
\(\epsilon_e\) の marker は \(s_E\) を選び、\(s_F\) と別の free scope を選ばない。
\(X\) が持つ他の source-owned key の通常 output discharge は、これとは別に行う。
一つの \texttt{'e?} を理由にそれらも一括で消すとは解釈しない。

\begin{theorem}[直接の compose の source owner と identity 出力]
\label{thm:compose-source}
上記 source 導出と group recipe に対し、\(g\,x\) から \(f\) の入力を通った一つの
同じ typed contribution \(w\) が、\(X\) の印付き output へ実際に到達するとする。
\(s_E\) の receiver はその時点で active な \(r_X\) であり、RC はその target の
\(s_E\) を解除する。\(f=id\) のときも output の \(\epsilon_f\) に共有される同じ寄与へ作用する。
各 effect/instance/attribution と、\(s_E\) 以外の保護はその release で変更されない。
内側で消費され到達しない最初の寄与には release は起こらない。
\end{theorem}
\begin{proof}
補題\ref{lem:compose-static}と EnterGroup から、\(s_E\) の起点は実際の
\(X\) entry にあり、元 receiver は \(r_X\)。完成した scope 表を別に仮定していない。
(C-Output) の release は \(\Scope_{r_X}\) の内側であるため、
その frame に到達した request を処理する時点では \(r_X\) はまだ現在の context にある。
\(f,g\) の内部 Scope が先に切り出されることは \(r_X\) を除かない。

実際に \(g\) の出力から \(f\) の carrier を通った \(w\) は、
同じ source port から O-Intro/transport の route と input stamp \(u_f\) を持つ。
従って From と Target は補題\ref{lem:inventory}、\ref{lem:target}により構成される。
\(g\) の genuine body route では \(u_g\notin T\) なので Protect が \(s_E\) を加える。
\(g\) の入力由来など、既に \(u_g\in T\) の場合は、その scope を新規に追加したとは主張しない。
どちらの場合も \(s_E\) の解除は同じ scope だけに作用する。
\(f\) へ入る route の Stamp により \(u_f\in T\)。
identity の source body はその同じ入力値/view を返すため、この tag と attribution が保たれる。
その \(f\) 出力の Protect は新しい \(s_F\) を追加せず、既存の \(s_F\) はそのまま保持する。
identity の source 名による特別規則はない。

従って RC の全条件が揃い、\((s_E,t)\) を \(L\) に追加する。
RP は whole output の各 witness を query 前に読むので、同じ origin が \(\epsilon_f\) の側に
存在しても未解除のコピーを残さない。各元 witness の画像を一度だけ構成し、
解除前の画像を同じ target の別 alternative として追加する規則はない。
\(s_E\) と異なる scope の bit、operation/attribution/Grant はこの transfer が変更しない。

handler が最初の request を消費すれば release frame の Req case は実行されない。
Value entry の場合は receipt の後に g を実行し、Computation entry の場合は実際の force がなければ
この \(w\) 自体が発生しないので、使用しない引数の効果を seed しない。
二回の通常の \(X\) 呼出しは fresh \(r_X,s_E,t\) を持つため、一方の release は他方を変えない。
raw resume、normal exit 後の latent 値は補題\ref{lem:life}の元 receiver 条件をそのまま使う。
\end{proof}

この定理は\emph{source key とその実際の entry/output}を供給する点で、
旧稿の「指定済み \(s_e\) と live receiver があれば」という局所計算より強い。
一方、これを全 source の導入算法の完全性、全 provider の将来の effect membership、
表示した scheme の主型性や native の全 route の再現と同じにしてはならない。
上述の conditional recipe を全表層言語へ採用・一般化する根拠は、依然として別の対応である。

\paragraph{一般化を妨げる既知の source 反例。}
\texttt{my direct f x = f x} と
\texttt{my wrapped f x = (\{ my unused = 0; f \}) x} は、元 formal と named group の情報を共有する。
しかし frozen Oracle の local-callee guard は実際の \(\mathsf{Var}\) だけであり、
後者の \(\mathsf{Block}\) を通さない。
そのため、本節の constructor 対応を「最後に同じ formal を返す全 source」へ無条件に広げられない。
これを言語の新しい拒否条件にして証明を閉じることもできない。
2026年10月15日の forwarding 決定は、ordinary forwarding と receiver-relative な保護を要求し、
具体的な suspension/escape/reentry 表現までは選んでいない。
本節の一般化では、alias/返却/capture の authentic source origin を、
元の source receiving boundary と一貫して結ぶ構成が次の義務になる。

\section{捕捉範囲をソースから構成する}
\label{sec:capture}

\subsection{既存の条件付き構成から、どの前提を取り除くか}

\nolinkurl{notes/theory/2026-10-15-effect-scope-ownership-capture.md} は、
与えられた region forest と capture frontier を単射で複製する局所補題を持つ。
しかし frontier の source 生成と、将来現れる scope を含む inventory の完全性は前提であった。
本節では source の実行意味論に限り、この二つを外部入力にしない構成を与える。
推論器が保存する制約 graph の FreshUse は別の対象である。
そこへ本節を適用するには、制約 graph と以下の closure/environment の対応が依然として必要になる。

まず次の二種類を構文上区別する。
\[
 \underbrace{r,b,s,t,o}_{\text{既に確保した証拠定数}}
 \qquad\hbox{と}\qquad
 \underbrace{\nu z.\,c}_{\text{まだ実行していない確保命令の binder}}.
 \tag{C-Sorts}
\]
同じ source label を二度使うこと、scheme の AllE、閉包の複写は左の定数を freshen しない。
右の \(z\) は code 中の局所束縛変数であり、capture-avoiding な alpha 変換を受ける。
\(\nu\) を実際に実行したときに限り、全既存 support と異なる定数が左へ追加される。
source に \(\nu\) 記法を追加するのではなく、source entry から生成する命令をこう表示する。

\subsection{有限な recipe environment と閉包}

各 source Lambda/delay/handler arm と生成 adapter に code label \(\ell\) を付ける。
\(\ell\) の自由参照を三つの sort に分ける。
\[
 F_v(\ell),\quad F_q(\ell),\quad F_\omega(\ell).
 \tag{C-Free}
\]
順に項の binder、qualified 証拠 binder、recipe が読む owner/origin/target の binder である。
\(F_\omega\) は実行済みの定数集合ではない。source の lexical binder を指す有限の変数集合である。
型/効果に残る自由変数も、全 packet に属する型環境として別に保持する。

構成手順は次のとおり固定する。
\begin{enumerate}
\item source の有限導出を走査して、各 primitive command が読む変数と、
その命令が束縛する変数を記録する。通常 entry の receiver、formal の origin、group slot、
marker target はその entry/output の局所変数である。
\item \(\mathsf{FV}(c_1;c_2)=\mathsf{FV}(c_1)\cup
(\mathsf{FV}(c_2)\setminus\mathsf{bind}(c_1))\)、
\(\mathsf{FV}(\nu z.c)=\mathsf{FV}(c)\setminus\{z\}\) として再帰する。
構造値の子、annotation template、adapter、dictionary field にも同じ規則を使う。
\item 入れ子の closure を作る命令は、その closure の自由参照を読む命令でもある。
従って孫の code が使う外側の owner も、途中の closure が保持する。
遅延された実行は自由参照の解析から削除しない。
\item 再帰 code label は展開しない。有限な label/lexical-binder 対について、
上の依存集合の最小固定点を有限 powerset 上で取る。
各増加で新しい対が一つ以上加わるので、高々その対の総数回の増加で停止する。
局所で束縛される変数は、その scope を越える前に除く。
\end{enumerate}
これは型の graph の全 latent shape を走査して保護を作る算法ではない。
\(F_\omega\) は既に決めた source recipe の\emph{自由参照}だけである。
次節の候補政策で必要になる lexical-origin 表も、その命令の明示した入力として数える。

閉包の具体値は
\[
 \mathsf{clo}(\ell,\eta_v|_{F_v(\ell)},
                    \eta_q|_{F_q(\ell)},\eta_\omega|_{F_\omega(\ell)})
 \tag{C-Close}
\]
とする。各環境 entry は既存の値/view または証拠への参照である。
source の guarded 値 SCC は、先に有限のセル名と signature を確保し、
closure/構造値だけを登録してから公開する。
code が指す再帰セルにはこの同じ有限表を使う。computed root を未初期化のまま読まない。
この処理は source 値 SCC の既存の前提を実現するだけで、全 linker の対応を新たに証明したものではない。

\begin{lemma}[source からの環境充足]
\label{lem:closure-env}
第\ref{sec:all-source}節の有限 source 導出と、自由参照を明示した有限 recipe から、
C-Free/C-Close を構成できる。実行中に closure の code が自由参照を読むとき、
その参照は捕捉環境、現在の formal、または既に実行した局所 binder に存在する。
完成済みの capture frontier、動的 scope 一覧、将来の実行履歴を入力として要しない。
\end{lemma}
\begin{proof}
構成の停止は有限対の増加回数から従う。
code の変数使用とその lexical binder を結ぶ経路に帰納する。
同じ scope の使用は現在環境にある。closure の境を越える自由使用は C-Free の union に入り、
C-Close の制限写像でその値が保存される。
二つ以上の境を越える場合は、内側 closure を作る命令自身の自由参照に同じ変数が入るため、
外側の各捕捉に帰納仮説を繰り返せる。
局所 binder は使用より前の実際の entry/let で値を得る。
再帰参照は登録済みセルであり、body の再帰的展開を証明に用いない。
pattern 成功環境はその成功枝だけを渡し、Or では共通 binder だけを渡すので、
失敗枝専用の変数をこの解析が正当化することもない。
raw continuation の saved environment には同じ判断を保持する。
source の全式/pattern の残りは子の自由参照の union と binder 除去で尽くされる。
\end{proof}

\subsection{値経路の生成を全構文で固定する}

recipe の自由参照を捕捉することと、値の由来を外へ投射することは異なる。
例えば \(\lambda z.p\) の環境には \(p\) があるが、closure 値そのものを
\(p\) と同じ callable として扱う規則はない。
ここを混同すると、捕捉しただけの callback の保護を新しい関数の body effect に付けてしまう。

\textbf{二つの問い合わせを区別する。}
既存の boundary evidence が値の signature path に沿って運ばれることと、
次節で新しい omission profile を作るために「この callable は元 formal のどれか」を調べることは別である。
後者を root-origin query と呼ぶ。
新しい closure root に捕捉値の formal identity を投射しない規則は、
前者の既存 evidence を消す規則ではない。
実際の Call の result boundary や構造 adapter が signature の対応した latent path に
既存 boundary reference を運ぶ場合には、その View/Flow は従来どおり保持する。
その path に exposure が存在しなければ追加しない。
「返した値だから全ての深さにコピーする」ことも、「新しい closure だから全て落とす」こともない。

値/view ごとに、その値の\emph{当該 component}へ至る actual-route descriptor を保持する。
値の取得に対応する新しい immutable view node から、既存 node へ後方参照を置く。
edge には typed path の prefix/除去と、実際に選んだ child を記す。
型付き graph の全ての到達可能な値を同じ root provenance に union しない。
生成規則は次の表で定義する。

\begin{longtable}{>{\raggedright\arraybackslash}p{.25\textwidth}>{\raggedright\arraybackslash}p{.67\textwidth}}
\toprule source 操作 & 実際に追加する経路と保持するもの \\
\midrule\endhead
Var、値 let、Block の結果 & lookup/成功 binder/最後の値への一辺。prefix の計算が返した別の値の由来は足さない。AllE は同じ値参照を使う。\\
Lambda、delay & 新しい callable/carrier の root は新しい introduction。自由値は C-Close により捕捉する。捕捉値の由来はその root へ投射しない。呼出し/force 後は実際の body 結果の経路を使う。\\
Tuple、Record、variant、constructor & 実際の子へ component prefix を付けて保持。record spread は評価済み可視 map の選ばれた field だけ。上書きされた field の別経路を同じ出力 field に追加しない。\\
field/payload projection、構造 pattern & 一致する component prefix を除く。入力全体に付いた source profile は、その許された typed-path derivative として読む。無関係な sibling は含めない。\\
As、Or、Case & As の元入力と子 binder を区別する。Or/Case は実際の成功枝とその直接境界。失敗枝の値由来を成功枝へ union しない。失敗までの実行効果は取り消さない。\\
default、guard & 実際に実行した計算の値/効果だけ。default に到達しなければその値経路を作らない。\\
App & callee の取得、既存の receipt/argument carrier、実際の body の戻り値と result boundary を保持。signature の実際の対応 path の evidence を運ぶ。型一致だけで引数の formal identity を結果へ足さない。identity body の Var は通常の lookup としてその値を返す。\\
RefSet、reference read & write は実際の RHS view を保存する。read は現在格納されている view から新しい取得辺を作る。reference の address equality だけから任意の過去の内容を復元しない。\\
Catch、raw resume & 選ばれた arm/return の値経路と saved context の同じ環境。resume の応答値は本当に受け取った値。saved evidence は恒等。\\
Select、qualified evidence & field はその field、method は構成した closure、dictionary は選ばれた evidence field。行政的な eapp は普通の source entry でも新しい由来の種でもない。\\
PrimitiveOp、direct adapter & 固定された primitive/adapter 規則が返す子または新しい値。user cast は普通の App として body の結果を使う。型等式だけから由来の辺を補わない。\\
\bottomrule
\end{longtable}

古い view node に新しい親 edge を書き足さない。
再帰値や共有セルの object graph が循環していても、取得の度に作る view-occurrence の辺は
その時点で既に存在する node へ向かう。query は有限の path derivative と visited 対を使う。
object graph 全体の再帰的な展開や、未来の write の先取りは行わない。
将来の read が新しい内容を見る場合、その時点で新しい取得 node が作られる。

\begin{lemma}[実行に沿った由来の正確性]
\label{lem:source-route}
上の構成による任意の有限 prefix について、query の各答えは、実行された entry から
当該値 component への一つの actual typed route を持つ。
逆に、この表の操作で実際に作られた route は query に現れる。
未実行の body、失敗枝の値、同型の別の値は答えを増やさない。
\end{lemma}
\begin{proof}
各新しい view node の作成順に帰納する。
entry は一つの新しい起点、lookup/return は既存の経路への一辺、
構造 introduction は実際の child への prefix、projection は一致する prefix の除去である。
選択枝と record の可視 field は一つに確定しているので、不実行の sibling を逆像に含めない。
closure の root と captured environment を分けたため、環境の存在だけで root への辺は生じない。
実際の body がその環境値を返す場合は lookup/return の場合で得られる。
read はその時点の stored view を一つ指すため、過去の write との架空の union もない。
逆向きは route の最後の生成規則で場合分けし、その前半に帰納仮説を適用する。
finite prefix の object/view 表は有限であり、visited 対による query は停止する。
pattern request はまだ成功値を返していないため、値の経路を早く生成しない。
以上が表の全生成 case である。
\end{proof}

ここでいう完全性は actual-route query の完全性であり、未知の全実行を静的に正確に予測する
算法の完全性ではない。両方を同じ「由来解析」と呼んで置き換えない。

\subsection{捕捉・返却・一般化で過去の証拠を改名しない}

有限な実行状態 \(H\) は値セル、view node、closure 環境、Boundary、target、saved context の
有限表を持つ。終了した receiver の identity/由来表は残し、activity だけを落とす。
本意味論は garbage collection を規定しない。
従って使われなくなった metadata が残ることは、未定義の参照を読む口実にならない。
query の対象集合は、必要な entry からこの有限表を辿ってその時点で構成する。
sealed な将来 inventory や、複写の度の region-locality 分類を必要としない。

\begin{theorem}[source から構成した捕捉の保存]
\label{thm:capture}
上の closure と actual-route 構成、および \(\Core_{\rm life}\) の下では、
alias、構造値、返された関数、共有セル、rank-2 の複数使用、再帰と複数回の raw resume を含む
任意の有限実行 prefix について、次が成り立つ。
\begin{enumerate}
\item code の自由参照は充足され、既存の \(b,r,s,t\) とそれらの equality は保存される。
\item 同じ値を二度使うことはその値に属する既存 scope の複製ではない。
各使用で到達した未実行の通常 entry だけが fresh receiver/scope を作る。
\item 捕捉・返却・AllE 自体は Observe、Grant、release、receiver の再活性化を行わない。
元 receiver が失効した保護は、後の closure 使用だけでは復活しない。
\item 新しい普通の呼出しと、保存された entry 後の suffix の再開を区別する。
再開前に確保した evidence は固定され、再開後に本当に到達した新しい Call は fresh entry を行う。
\end{enumerate}
\end{theorem}
\begin{proof}
状態の well-formedness を、全参照が表の同じ sort の entry を指すこと、
各 closure が C-Free を満たすこと、全 scope が一つの確保済み Boundary に遡ることと定義する。
初期状態は有限な宣言/値表の構成からこれを満たす。
closure 作成は補題\ref{lem:closure-env}で必要な参照を制限写像に保存し、新しい証拠定数を作らない。
alias/構造/返却/read は補題\ref{lem:source-route}の実際の辺を足すだけである。
write は既に well-formed な RHS view を保存するので、循環した object graph でも参照は定義済みである。

AllI/AllE では、型と未実行 recipe の binder に同じ capture-avoiding な代入を施す。
動的 evidence 定数はその代入の定義域にない。
同じ closure を二つの型 instance で使っても、捕捉環境は同じ参照を持つ。
\(\nu z\) に実際に到達した場合は、有限 support の外に fresh 名を一つ取る。
この延長で全旧参照を固定するため equality は保たれ、新しい entry との混同もない。
再帰は同じ有限 code label の実行であり、呼出しごとの環境/確保について今の議論を繰り返す。

handler の捕捉は frame/environment の保存であり、その保存時に新たな View を実行しない。
raw resume は補題\ref{lem:life}の borrow/own と旧 evidence の恒等写像を使う。
既に通過した entry 命令は saved suffix の前にあるので再実行されず、
suffix 内の後続 Call は普通の entry の場合に含まれる。
leave は activity を減らすだけで表の参照を消さず、捕捉/返却に activity を足す規則はない。
以上の一段保存を prefix 長について繰り返す。
この帰納は停止する実行だけを仮定せず、発散する実行の全有限 prefix にも適用できる。
\end{proof}

この定理で frontier と将来 inventory の外部仮定を除けたのは、
\emph{実行中の既存値/evidence を複製しない source 意味論}においてである。
制約 graph の generic 部分を複製する推論器については、どれが code 内の binder に対応し、
どれが捕捉済み定数に対応するかの証明が必要である。
既存の region-clone 補題が不要になった、または native FreshUse がこの定理を既に満たす、とは結論しない。

\paragraph{寿命を区別する三つの導出。}
\begin{enumerate}
\item receiver \(r\) の実行中に callback を含む closure \(k\) を作り、同じ実行中に \(k\) を呼ぶ。
C-Close は \(r\) を固定し、\(\Act(r)\) は真なので、既存 protection は通常 query に残る。
closure を作っただけではその callback の event は発生しない。
\item 同じ \(k\) を返して \(r\) の正常完了後に呼ぶ。
保存された参照はなお \(r\)、新しい \(k\) の実行 owner は fresh \(r'\ne r\) である。
従って古い protection は inactive である。
新しい呼出しで本当に受け取った引数への別の protection は、この結論と両立する。
\item 同じ partial application を二度使い、末尾の通常 entry でそれぞれ \(r_1,r_2\) を得る。
早い引数の captured view は同じ参照だが、末尾の receipt は別の確保命令の実行なので、
\(r_1\ne r_2\)、\(b_1\ne b_2\)、それらの scope/target も別である。
一方の target への RC は他方の \(L\) entry を作らない。
\end{enumerate}
各項は C-Close の参照の恒等性、通常 entry の fresh 条件、RC の pair equality を順に適用した導出である。
「返却」と「別の通常 entry」と「raw resume」を同じ fresh-use 操作にすると、この三つを区別できない。

\section{non-Var の由来と owner を決める全域な候補規則}
\label{sec:flow-policy}

前節は「与えられた recipe を正しく捕捉する」部分を閉じた。
この節では recipe 自身の source 生成を一つ具体的に定める。
政策を \(\Pi_{\rm flow}\) と呼ぶ。direct named case 以外への拡張は\emph{未採用の候補}である。
source 構文を拒否せず、Var-only guard を一般化する場所を明示する。
元の receiver を常に current owner に置き換える規則ではない。

\subsection{元 formal の静的 identity と、実際の binding token}

無注釈 formal \(p\) の実際の binding を公開する際に fresh token \(o\) を作り、
\[
  \mathsf{Seed}(o)=(p,b_o,r_o)
  \tag{F-Seed}
\]
を記録する。\(b_o,r_o\) はその\emph{実際の receipt} の Boundary/receiver である。
token 自体は effect、Observe、Admit/Grant を作らない。
omission profile は実際の upper-use で露出する output port にのみ適用される。
明示 callback wildcard、既知外部 Name、ただの local let に F-Seed はない。
明示 concrete/computation annotation は既存の source profile を通常 receipt で使い、
無注釈 origin へ変換しない。

receipt が force/pattern より先である順序は維持する。
entry で \(b_o,r_o\) と slot を予約し、formal の値/token 環境への登録は
その binder を作る subpattern の成功後に行う。
Value formal の入力計算が request を出す段階で、未確定な binder の値を読まない。
Computation formal の carrier も通常の保持/force 規則を使い、F-Seed のために先に force しない。
複合 parameter pattern の場合、\(p\) はその成功環境へ公開する解決済み binder ID とし、
元 parameter slot と実際の component path を一緒に持つ。
Or 内部の default 等には、その試行で既に成功した branch-local binder の環境を渡す。
外への公開 binder は共通 merge binder に対して別に作り、選択された枝の直接境界で値を埋める。
不一致なら branch-local 環境を捨てるが、その試行の effect/共有状態は巻き戻さない。
As の whole binder と component binder は別の source binder として記録する。
それぞれの annotation 有無は元の source から読み、無注釈でない子を omission seed にしない。
これらは同じ実際の parameter receipt を参照し、path/binder の違いを \(j\) に残す。
raw resume がまだ行っていない binding 公開に到達すれば token は新たに作るが、
その起点の \(b_o,r_o\) は保存された元 receipt のままである。

formal を束縛した lexical 環境には \(p\mapsto(o,\chi_p)\) を保存する。
\(\chi_p\) はその source binding が現在参照する receiving slot であり、初めは \(b_o\) である。
token と slot の identity は別物である。
source の upper-use occurrence \(u\) には、その lexical scope にある formal の有限表
\(\Lambda(u)\) を付ける。
closure ではこの表が読む token 変数も \(F_\omega\) に入り、捕捉時の値を保持する。
callback の返却先で呼出し stack を検索して、同名の別 binding を探す規則はない。

formal の型付き root に token の由来を付ける。実際の alias/projection/adapter/body return は
前節の表でそれを運ぶ。root 全体の source profile の descendant は、実際に選んだ typed path の
derivative として遅延して読む。closure の捕捉環境を root の token 集合へ潰さない。
次の記号はこの query の答えであり、任意の型判断の追加前提ではない。
\[
 \mathsf{RootOrigin}(v,o,j,\tau).
 \tag{F-Root}
\]
\(j\) は元の profile port/path、\(\tau\) は実際の route である。
新しい source introduction が単に同じ型を持つことは、この関係の生成規則ではない。

\paragraph{upper-use の producer。}
upper-use は App の callee だけに限らない。
有限 source 導出の App 規則が作る complete Function demand、および、
引数/結果/注釈/field の\emph{実際の直接境界導出}が期待側に Function を露出する node を列挙する。
各 node に元の source rule と方向を付け、型の provider 側から逆向きに列挙しない。
入れ子の Function 引数では向きが反転するため、構成された比較の向きを使う。
文字列上の「期待型の中にある」という位置だけでは upper と判定しない。
direct boundary の child は有限な証拠木/DAG の既存 label を使う。
再帰型の無限展開や、既存 lower Function の全走査は不要である。
既存の値をその upper interface へ運ぶ時、\(\mathsf{Expose}_u\) 命令を実行する。
この命令は以下の F-Owner を使って consumer view の output profile を構成する。
profile が latent callable に属するなら、その callable の\emph{実際の実行時}まで Observe を行わない。
その場合の input tag も Expose 時には確保しない。
profile は \(s\) を自由定数として捕捉し、各実際の invocation に
\(\nu i.\,(\mathsf{Stamp}_i,\mathsf{Protect}_{s,i})\) を実行する template を保持する。
同じ view を二度呼んだ場合の tag は相異なる。
tag の確保はその既存 Call/force の metadata 処理であり、新しい source Call や receiver を挿入しない。
これにより、第一回の入力由来を第二回の入力由来と誤認することがない。
例えば callback を別の関数の Function 型引数へ渡す場合には、その実際の引数境界が producer になる。
そこで既に生成した profile は、後の callback 呼出しで改めて owner を選び直すことなく運ばれる。
これは本候補の宣言的 producer 規則であり、推論器が全 source upper constraint を列挙するという
未証明の実装命題を前提にしていない。

\subsection{source owner slot の構成}

named header の最後の Lambda を group terminal とする。
anonymous Lambda はその一つの formal の group として扱うのが本候補の追加規則である。
zero-parameter initializer、Block、Group、型適用は新しい group ではない。
source upper-use の innermost lexical function group を \(G(u)\) と書く。
function の外では \(G(u)=\bot\)、\(P_\bot=\varnothing\) とする。

named group の通常の terminal entry で、その group の無注釈 formal \(p\) ごとに
\(B_{G,p}\) を予約する。これは第\ref{sec:named}節の EnterGroup と同じ source entry であり、
現在の \(r_G\) と元 formal の source descriptor を記録する。
複合 parameter では元の有限 pattern 導出に現れる binder を列挙し、
まだ成功していない binder の slot は予約だけ行う。対応する成功時に token/view を埋める。
失敗枝の slot は外の環境に公開されず、単なる予約を Observe に使う規則はない。
ここで「埋める」は現在の lexical 環境を immutable に拡張することである。
既存の Boundary レコードや saved environment を上書きしない。
その group 自身の各 \(p\) の新しい lexical 環境を
\(p\mapsto(o,B_{G,p})\) とする。新しい body が捕捉する \(\chi_p\) はこの slot である。
ancestor の \(\chi_p\) は変更しない。
値環境には同じ raw value の新しい consumer view を置き、その entry の実際の受領辺を残す。
\(\mathsf{Seed}(o)\) の古い receiver を書き換えるわけではない。
新 scope の Origin は選んだ \(B_{G,p}\) の actual entry を指し、
token equality だけを境界の代用品にはしない。
一引数で original receipt と terminal が同じなら \(B_{G,p}=b_o\) とし、重複確保しない。
未使用の slot は output View に参照されないので、その存在だけでは保護 witness にならない。
profile の全 latent shape は展開せず、後の実際の port query が \(\mathsf{scopeKey}(B_{G,p},j)\) を読む。
予約と後の読み出しは、新たな source receiving event を二重に数える操作ではない。

実際の \(\mathsf{Expose}_u\) が受け取った値 \(v\) について、F-Root の各答えへ
以下を実行する。App の callee もこの場合の一つである。
以下の式では \(p=\mathsf{Seed}(o).p\)、\(\eta_\omega(p)=(o_p,\chi_p)\) とする。
\[
 \mathsf{Eligible}(u,o)\iff p\in\Lambda(u)\ \land\ o_p=o.
\]
\[
 \mathsf{OwnerSlot}(u,o)=
 \begin{cases}
 \chi_p,&\mathsf{Eligible}(u,o),\\
 \mathsf{none},&\neg\mathsf{Eligible}(u,o).
 \end{cases}
 \tag{F-Owner}
\]
同じ group の body では \(\chi_p=B_{G(u),p}\)。
nested closure では、その closure が捕捉した同じ \(\chi_p\) を使う。
none は通常 App/境界の拒否ではない。新しい omission profile を導入せず、既存の view を使う。
source signature に既にある annotation boundary の通常処理も保持する。
\(p\) は source の解決済み binder、\(o\) は\emph{現在 code が捕捉したその binding}である。
同じ \(p\) の再帰的な別 activation の token をこの equality で弾く。
古い token に付いた既存 evidence を削除するわけではない。

選ばれた slot の \(s=\mathsf{scopeKey}(b,\epsilon,j)\) を、当該 upper-use の
\emph{出力 View} の Protect にだけ入れる。
source の input carrier の実際の実行には N-Transfer の Stamp を使う。
App ではこの View が通常 Call を囲み、latent upper interface では実際の callable 実行を囲む。
従って内側 handler の判定より前に実際の emission から観測できる。
全 executing View は同じ actual contribution の input tag を参照し、
別の unstamped packet を勝手に選ばない。
lower provider の共有 graph、引数の型全体、既知外部関数へ逆方向の seed を追加しない。
複数 token が実際にある場合には全ての独立した slot を保持し、一つの代表へ統合しない。

\textbf{curried entry と raw resume。}
最初の引数を受け取って返す途中の closure は、その時点の \(p\mapsto o\) を捕捉する。
最後の引数の\emph{普通の呼出し}が、そこに fresh な terminal slot を加える。
同じ partial application を二度使うと、早い receipt は共有され、terminal slot は異なる。
raw resume は saved environment の \(B_{G,p}\) を固定する。
新しい実行 control owner があっても、F-Owner はそれを代用しない。
捕捉した ancestor の \(\chi_p\) の receiver が既に inactive なら、その保護も inactive のままである。
「近くの live owner が見つかるまで検索する」分岐は存在しない。

\paragraph{構成途中で棄却した owner 規則。}
ancestor なら常に最初の \(b_o\) へ戻す案は採らない。
これは凍結 Oracle の C2 の静的 introduced-frame 選択を、そのまま物理的な早い receipt に
対応させようとした案である。次の名前解決済み source 表示が違いを示す。
\begin{Verbatim}[fontsize=\small]
my direct cb x = catch cb x { ping(), _ -> 1; v -> v }
my nested cb x = {
    my step z = cb z;
    catch step x { ping(), _ -> 1; v -> v }
}
\end{Verbatim}
\(cb\) が \(\mathsf{ping}:1\to\mathsf{int}\) を放出し、明示の捕捉許可はない場合を考える。
curried な最初の \(cb\) entry は残りの関数を返して終了し、\(x\) entry の \(r_X\) は現在 live である。
direct の scope は \(B_{X,cb}\) を持つ。
棄却案の nested では、helper 内の上位 \(cb\) use が inactive な \(b_o\) だけを参照するため、
外側の local Catch が ping を処理できてしまう。
source 型が一致することだけではこの差は排除されない。
charter 第13節の「引数から捕捉への経路変更で無許可の捕捉を可能にしない」条件に対し、
この対応は不足する。

修正版では実際の \(X\) entry が \(\chi_{cb}=B_{X,cb}\) を作り、
step の C-Close がその同じ参照を保存する。
helper 内の actual upper-use も \(B_{X,cb}\) を読むので、live な \(r_X\) の protection が残り、
両方の local Catch はこの witness により遮断される。
外側の scope が終了すれば、step を返して保持していてもこの witness は inactive になる。
これは source 上の紙上判別であり、compiler 実行の観測ではない。
静的 C2 の値型構成を否定するものでもないが、その全 frame を元の物理 receipt に結びつける
source/runtime 対応は、本候補では使用できない。

\begin{lemma}[全域な recipe 構成と source 起点]
\label{lem:flow-formation}
有限 source 導出から \(\Pi_{\rm flow}\) の slot/環境/query 命令を有限に生成できる。
その任意の有限実行 prefix で、導入する omission scope は一つの actual source receipt または
terminal entry に遡り、元 receiver を変更しない。
\(\mathsf{ProtectedVarAt}\)、完成した origin/owner 表、capture frontier を自由な前提として要しない。
\end{lemma}
\begin{proof}
lexical 解決で各 \(p,G(u),\Lambda(u)\) は有限に決まる。
upper-use の列挙は App と有限な直接境界導出の指定 node への再帰である。
これらの node の実際の consumer に Expose を挿入するため、型の同形性から起点を捏造しない。
予約 slot は有限な group/formal 対、各 use の output profile はその source 型導出の port template である。
再帰 code は同じ label を指すので命令数を増やさない。
closure の token/slot 参照は C-Free により構成する。
entry で F-Seed/EnterGroup が実際の受領値と receiver を記録する。
RootOrigin の答えは補題\ref{lem:source-route}からその受領へ遡る。
Eligible の equality が actual binding と異なる再帰 instance を除く。
同じ group の場合は実際の terminal entry の slot、ancestor の場合は捕捉済み \(\chi_p\) を読む。
\(\chi_p\) は元 receipt または実際の source terminal entry でしか束縛されない。
どちらも既存の表を読む全域操作であり、missing origin は none の場合になる。
profile の適用は actual upper-use の output だけなので、型の未来の解を起点にする必要がない。
inactive な slot を読む場合も参照は定義済みであり、activity query が保護を無効とする。
全ての確保は上記二つの source entry で尽くされる。
\end{proof}

\begin{proposition}[別名と新しい関数の区別、および再帰 instance の区別]
\label{prop:alias-origin}
\(\Pi_{\rm flow}\) では次の source 導出が成り立つ。
\begin{enumerate}
\item \texttt{my direct f x = f x} と
\texttt{my wrapped f x = (\{ my unused = 0; f \}) x} は、
対応する実際の call で同じ元 \(f\) token と terminal slot を選ぶ。
\item \texttt{my a = f; a x} も同じである。
一方、\texttt{my a z = f z; a x} の \(a\) は新しい closure root であり、
外側の \(a\) 呼出しを \(f\) の呼出しに偽装しない。
内側 \(f\) 呼出しは、その closure が捕捉した receiving slot \(\chi_f\) を選ぶ。
\item 同じ source formal \(p\) に二つの actual binding token \(o_1\ne o_2\) があっても、
現在の \(\eta_\omega(p)=(o_2,\chi_p)\) を使う code は \(o_1\) を current terminal slot へ付け替えない。
\end{enumerate}
\end{proposition}
\begin{proof}
最初の Block は一つの literal binding の後に同じ \(f\) の Var を返す。
source recipe の自由表と group は変わらず、actual route が lookup 一辺だけ延長される。
RootOrigin の \(o\)、Eligible の equality、OwnerSlot の結果は同じである。
値 alias も同じ lookup/return の場合である。
新しい関数の場合は C-Close の root 規則により \(f\) token を \(a\) root に投射しない。
その body 内では、\(f\) の lexical な token と \(\chi_f\) を C-Close がそのまま供給する。
第三項では equality の偽からその lexical slot を選ぶ分岐を使えない。
旧 \(o_1\) の実際の経路と既存 protection は別に保持されるため、これは消去ではない。
\end{proof}

第1項は凍結 Oracle の C3 と異なる、候補規則自身の結論である。
原 source を Var の形に制限したり、旧 Oracle と一致すると主張したりしていない。
第2項の後半は eta 展開の全観測同値を主張しない。
新しい受領境界や annotation を挿入する変換が、全て無意味であるとは仮定できない。
multiargument outer の早い formal を nested helper が使う場合にも、
本候補は実際の outer terminal entry の slot を運ぶ。
古い early receipt を復活させる操作はなく、terminal entry がなければ新しい slot もない。
この構成で上の捕捉判別を満たしたことと、全 source の意味論的対応は区別する。

\begin{lemma}[lexical な receiving slot の捕捉保存]
\label{lem:home-capture}
source 環境が \(p\mapsto(o,\chi_p)\) を持つとする。
その環境から何段の nested closure を作り、alias/構造 field/返却を経由して実行しても、
その \(p\) を同じ lexical binding として使う upper exposure は同じ \(\chi_p\) を読む。
\(p\) の宣言 group の新しい普通の entry、明示 release、expiry はそれぞれ別の遷移として扱う。
\end{lemma}
\begin{proof}
closure を越える一段は C-Free が \(o,\chi_p\) の両方を自由参照に含め、
C-Close が同じ組を保存する。複数段はこの一段の帰納である。
alias/projection/返却は code の捕捉環境を書き換えない。
実行する内側 group は自分の formal の slot だけを束縛するので、ancestor の組を変更しない。
同じ binding token を持つ RootOrigin に対する F-Owner は保存された \(\chi_p\) を返す。
profile の実際の受領辺はその slot が作られた entry に遡る。
従って source の actual typed route を持つ同じ protection witness は捕捉だけでは失われない。
\(p\) の group の新しい entry なら新しい source 受領、明示 release なら RC/RP、
expiry なら activity の変更として、元の witness への作用を別々に判定できる。
raw resume はこれらの自由参照を変更しない。
\end{proof}

\subsection{direct compose と一般の返却値へ戻す}

\begin{theorem}[source 由来を持つ alias/捕捉/返却の合成]
\label{thm:flow-compose}
\(\Pi_{\rm flow}\) の source recipe は直接の named compose の
第\ref{sec:named}節の構成を再現する。
さらに値 alias、純粋な Block の結果、実際に成功した field/pattern の取得を通る
同じ formal への callee route に対して、対応する source binding が同じなら同じ slot を使う。
返された関数/多相 field/共有セルに運ばれた既存 evidence は、定理\ref{thm:capture}の寿命を保つ。
\end{theorem}
\begin{proof}
direct compose では両 callee の RootOrigin はその formal の token を含む。
両 Eligible が真、両 formal は現在 group のものなので、
slot は \(B_{X,g},B_{X,f}\)。これは補題\ref{lem:compose-static}の二つの key である。
予約した未使用 slot は実際の View に現れないため、存在だけで query の答えを追加しない。
各 alias/projection の場合、補題\ref{lem:source-route}が同じ token と実際の path derivative を返す。
同じ lexical binding の条件が F-Owner の同じ分岐を保証する。
Source の output marker はその slot の inventory を読むため、
定理\ref{thm:compose-source}の release と whole-output の計算をそのまま適用できる。
実際に返された closure/field/store の場合は C-Close と実際の値経路の生成で運ぶので、
定理\ref{thm:capture}から旧 receiver/target の同一性と expiry を得る。
新しい source Call があればその新しい receipt だけが別の scope を作る。
\end{proof}

これは、同じ source binding を持つ actual callee route についての一般化である。
その条件を「同じ型である」「何らかの closure 内に同じ callback がある」に弱めてはならない。
一般の関数の body がどの値を返すかを静的に当てる必要はなく、実際に返した値の経路を使う。
そのため未使用 computation や発散する callback を、由来を知るために先に実行することもない。

\section{ソースからの健全性定理を一つに合成する}
\label{sec:source-family}

\subsection{未選択の政策と、型健全性に不要な推論義務}

\(\Pi_{\rm flow}\) と、各 marker の有限な ingress 選択政策 \(A\) を固定した体系を
\(\Src_{\Pi_{\rm flow},A}\) と書く。
\(A\) は lexical binder と source ingress occurrence/template だけを選び、
異なる binder/instance を row equality で統合しないものに限る。
また、本稿の単一 ingress と直接 compose について既に固定した対応に一致させる。
その唯一の対象を空集合へ変える政策を、元の \texttt{'e?} の解釈としては含めない。
複数 ingress がある場合の \(A\) をユーザーが既に選んだ、とは仮定しない。
以下はこの明記した体系族についての一様な定理である。
選ばれていない政策を隠した「Yulang 全体」の一意な意味論とは呼ばない。

declared source typing の証人は、有限な値型/全称/境界/選択/role の規則木と、
各生成 port の有限 effect 集合である。
後者に対して、生成した全 mandatory edge、全共同入力、cut-context summary、
immediate/latent/reentry の inclusion を検査する。
型付けに effect inclusion の証拠を要求することは、推論器がそれを発見するという仮定ではない。
実行の安全性そのものや完成済み runtime metadata 表を証人として渡すこともしない。

\subsection{新しい source 処理の有限表示}

source の finite template/label/binder、token/Boundary/receiver の equality、
activity、既存の incidence/From の共同座標を complete cell に持たせる。
closure 環境を無限に inline 展開しない。
既存の finite interface が表さない hidden origin/route に依存する場合には、
\emph{整合する全 output cell}を許す関係にする。
この fallback は powerset の最大元を含む有限列挙であり、精度の主張はしない。
inclusion が満たせなければ、その小さい effect annotation の導出は得られない。
これを source の動的失敗や、隠れた正確な membership oracle に置き換えない。

新しい転送は三種類で尽くされる。
\begin{enumerate}
\item C-Close、alias、captured-environment の lookup は evidence の恒等コピーと
実際の typed route の追加。新しい View occurrence と既存 Boundary を別 sort にする。
\item F-Seed/terminal entry は既存の fresh/receipt の relation と、
source label/actual binding の有限な同時登録。未使用 slot に Observe を生成しない。
\item F-Owner は RootOrigin、lexical 表との token equality、保存された receiving slot の取得。
選択後は N-Transfer と旧 O-Target/RC/RP。unknown の場合には全整合分岐を残す。
\end{enumerate}
closure に保存する環境と code の自由参照の対応は補題\ref{lem:closure-env}が与える。
有限 interface が closure の全内部 graph を正確に表現するという追加仮定は置かない。

\begin{lemma}[source closure/owner transfer の表示保存]
\label{lem:flow-transfer}
上記有限 relation は、\(\Pi_{\rm flow}\) の各具体 metadata step の出力を含む。
全 graph edge は有限集合間の単調 relation image であり、固定した finite signature で検査できる。
\end{lemma}
\begin{proof}
copy/route/fresh の場合は具体操作と同じ既存 relation を使う。
F-Owner の既知座標の場合には、実際の token equality と slot 取得が列挙した一つの分岐になる。
未知の場合は全整合 output cell に具体出力の cell も含まれる。
scope/target の選択後には補題\ref{lem:group-transfer}と\ref{lem:target}を適用する。
どの relation も有限 signature の powerset 上にあり、入力集合を増やして既存の画像が消える操作はない。
従って単調である。これは効果の最小集合の精度や、推論器の探索完全性の証明ではない。
\end{proof}

\begin{theorem}[固定した source 政策族に対する推論器独立の健全性]
\label{thm:source-safety-family}
有限で well-formed な source 宣言表、閉じた名前解決済み source \(e\)、
固定した \(\Pi_{\rm flow},A\) の下で上記の宣言的導出
\(\vdash_s e:A_0!\Phi\) が与えられたとする。
その導出から実行 code、closure 環境構成、source origin/owner/target 命令、
有限な effect 検査 graph を構成できる。
全有限実行 prefix で以下が成り立つ。
\begin{enumerate}
\item 値型が保存され、値、型に合う request、列挙した failure、次の遷移のいずれかである。
\(\Bad\) へ到達しない。
\item outward request と latent/reentry の実行は \(\Phi\) に含まれる。
空の outward 契約から request は出ない。
\item protection は actual source entry と actual typed route に遡り、
handler は同じ共同 packet の eligibility を満たす場合だけ request を処理する。
単なる alias、捕捉、型 instance、raw resume で旧 receiver は復活しない。
\item release は選択された origin と元 receiver の実際の crossing に限り、
その whole-output view に作用する。独立 scope、effect membership、attribution、Grant は保つ。
\end{enumerate}
この定理の前提に推論器、resolver の停止性、主型性、完成した capture/receiver 表はない。
\end{theorem}
\begin{proof}
定理\ref{thm:source-value}と補題\ref{lem:source-poly}が source 規則木から実行型を作る。
選択の各使用は有限の BuildSelect 導出を組み込み、各 body は rigid な宣言検査済みである。
補題\ref{lem:flow-formation}が source から entry/owner 命令を作り、
補題\ref{lem:closure-env}がその自由参照の環境を構成する。
補題\ref{lem:source-route}と\ref{lem:target}が各 actual origin と選択政策 \(A\) から
target を作る。従って完成した任意の metadata 表を型付けの前提に戻していない。

実行型については新 metadata が underlying value/code の型を変えず、
定理\ref{thm:life-value}の一段保存と canonical forms を使う。
型変数/qualified 証拠の代入は whole packet に同時に行い、動的定数は固定する。
型抽象・closure 環境の保存は定理\ref{thm:capture}により同じ型付き参照を保持する。

effect については生成 graph の全 inclusion を前提の有限証人で検査する。
各一段の具体 metadata を補題\ref{lem:flow-transfer}が表示し、
既存の seven-frame cut の各 typed hole にこの表示を入れる。
新しい scope delimiter の種類はなく、source lookup/予約は frame 内の有限処理である。
handler の No 試行後には current state を再検査し、raw continuation の suffix には
旧 evidence 恒等/現在 control の cut summary を使う。
定理\ref{thm:life-full}と\ref{thm:group-safety}の prefix 帰納が、全 outward/latent/reentry を与える。
unknown branch で effect 集合を広げた場合も、具体出力を含むので同じ保存が使える。

最後に origin の一意な起点は補題\ref{lem:flow-formation}、
captured evidence の寿命は定理\ref{thm:capture}、captured receiving slot は補題\ref{lem:home-capture}、
actual route は補題\ref{lem:source-route}による。
RC は元 receiver と現在の crossing を検査し、RP は同じ target の全 output witness を投射する。
これらの式は他の scope/attribution/Grant を変更しない。
従って項目3、4を各 prefix について同時に保持できる。
全証明は有限導出と有限 prefix の帰納であり、実行や推論の全域停止を要しない。
\end{proof}

\paragraph{この定理で完了したものと、完了していないもの。}
候補 \(\Pi_{\rm flow}\) では、13式/11 pattern の source recipe 生成と捕捉を
完成した runtime 表の外部仮定から取り出し、source 型健全性へ戻した。
ただし、候補の nested/anonymous owner 選択が全ての承認済み・意図された source 振舞いを満たすこと、
複数 ingress の \(A\) の選択、全表層構文からこの source への展開はこの定理の結論ではない。
型健全性は、別の型に合う handler を選んでしまう意味論の相違を自動的には反駁しない。
したがって「唯一の Yulang 意味論が閉じた」「独立査読だけが残る」という結論にはしない。

\section{dot selection との接続}
\label{sec:selection}

選択稿の有限 DeclOK、M/R/F、direct boundary、qualified role 証拠を用いる。
新しい source 全称規則が必要になる箇所を明示する。
\begin{enumerate}
\item \textbf{field。} record/nominal の宣言 field が \(\forall a.B\) なら、
field projection の source 結果はその全称型である。
使用式の AllE を \emph{projection の後}に置く。フィールド取得時に一つの \(T\) を選ばない。
\item \textbf{method。} member scheme の外側の binder と全 qualified 証拠は、
一つの use に一つの同時 instance を使う。
payload/result 内にある別の \(\forall\) はそのまま残し、使用時にだけ AllE で開く。
\item \textbf{impl。} body を rigid head と同じ candidate-wide 関連型の下で検査する。
多相 field は AllI と値制限を満たす body から作る。
既存の単相 body が一つの instance に合うことだけでは全称 field にならない。
\item \textbf{eabs/eapp。} 管理上の証拠適用は source Call を追加しない。
free な evidence の receiver は固定する。
一般の eabs の body にある source call は、その body を実際に実行するときだけ
普通の fresh receiver を導入する。純粋な dictionary 取得に release 状態変更を混ぜない。
\end{enumerate}

\begin{theorem}[source の多相選択から修正版への構成]
\label{thm:selection-source}
上記の有限 source 導出と DeclOK/選択/境界/role 導出から、
\(\Core_{\rm life}\) の型付き実行計画を構成できる。
field の内側の全称型、method の whole instance、captured receiver の同一性を保つ。
生成した計画が修正版の全 finite effect inclusion を満たせば、定理\ref{thm:life-full}の結論を得る。
\end{theorem}
\begin{proof}
定理\ref{thm:source-value}の Select case へ、選択稿の BuildSelect の構造帰納を入れる。
field は同じ payload 型の projection、method は rigid 検査済み body の一つの qualified instance、
role は有限証拠 DAG の child から作った dictionary field を使う。
全称導入/除去には補題\ref{lem:source-poly}を適用する。
eabs/eapp の substitution は値/型/証拠の同時置換であり、過去の receiver を freshen しない。
source Call だけが実際の receipt を作ることから、管理適用で余分な receiver を増やさない。
修正後の graph を生成して全 mandatory inclusion を検査した場合に限り、
定理\ref{thm:life-full}を適用する。旧 graph の小さい residual を前提なく再利用していない。
\end{proof}

group recipe を使う source では、同じ構成を \(\Core_{\rm group}\) に対して行い、
Stamp/Protect/EnterGroup の edge を生成してから定理\ref{thm:group-safety}を使う。
dictionary の取得で実行 owner や group recipe を発火させないため、
field から取り出した全称値の各 use と、実際の source Call の fresh entry が区別される。
これは source の rank-2、選択、group recipe を別々の不整合な核へ翻訳しないための合成である。
\(\Pi_{\rm flow}\) を使う場合には、取得した field/dictionary 値にも
補題\ref{lem:source-route}の actual component 経路を使い、
method closure の自由参照を C-Free から構成する。
従って定理\ref{thm:source-safety-family}の Select case に同じ BuildSelect を代入できる。
method の自由変数を捕捉したことだけで、その callback の origin を method root へ投射しない。

\begin{proposition}[選択の整合性と型健全性の分離]
\label{prop:selection-policy}
各 selection site で同じ宣言的結果型/効果契約を満たす有限選択証拠を一つ固定する。
どの証拠を固定しても、生成 code には定理\ref{thm:source-safety-family}が適用できる。
このため、resolver の全到達 snapshot で選択が一致するという命題は、この型健全性の前提ではない。
ただし、異なる証拠の生成 code が観測同値であるという結論は得られない。
\end{proposition}
\begin{proof}
BuildSelect は各証拠について同じ site の型/効果契約の導出を構成する。
それを該当する型付き穴に入れ、全 mandatory inclusion を検査する。
定理\ref{thm:source-safety-family}の証明は、その証拠をどの探索器が選んだかを参照しない。
他方、同じ \(1\to\mathsf{int}\) 型の二つの body が0と1を返すことは可能であり、
型/効果契約だけから観測同値は導けない。
\end{proof}

\paragraph{選択の一意性について何が残るか。}
有限な snapshot と走査順序が完全に与えられ、各検査の有限な真偽証拠があるとき、
first-success と順序付き fallback は高々一つの結果を返す。
証明は候補列の長さへの帰納であり、先頭が成功ならその identity、失敗なら残りの列へ進む。
これは resolver の政策に従った\emph{固定 snapshot の決定性}である。
別の到達可能 snapshot でも同じ候補が選ばれること、compact/raw head の一致、
役割実装の全体的一意性、SCC の Ready 公開の正当性は、この帰納法では得られない。
同じ型の別 body を選んでも値型が保存される可能性があるため、型健全性だけで穴を埋められない。

\section{親定理への適用と残る境界}
\label{sec:remaining}

今回の source からの構成は次を実際に減らした。
\begin{itemize}
\item 多相 payload と明示 instance の部分では、核 code を与える前提を、
source 宣言・rigid 検査・使用箇所の AllE の有限導出へ戻した。
\item receiver を過去の証拠ごと改名する source 対応を反駁し、
元の receiver と実行 owner を分けた修正版の保存・進行と finite transfer を与えた。
\item 実行された source 境界からの origin inventory と target 構成を与え、
任意の完成した scope/receiver 表を premise とする部分を縮めた。
\item compose の実際の Call/Scope 順を追い、元 receiver が live であるという
局所補題の前提を \(f\) の入口 receipt からは供給できないことを特定した。
既存の source constructor を適用して \(X\) の実際の entry を起点に置き換え、
直接の compose の source-to-scope/target と identity 出力まで紙上で接続した。
\item C-Free/C-Close と actual-route の生成により、実行意味論の capture frontier と
将来 inventory を外部から与える前提を除いた。alias/返却/共有/再帰/raw resume の保存を
定理\ref{thm:capture}へまとめた。
\item non-Var upper-use の政策 \(\Pi_{\rm flow}\) を一つ明示し、その source producer から
健全性定理\ref{thm:source-safety-family}まで合成した。
政策の採用・全期待挙動との一致と、政策を固定したときの安全性を区別した。
\end{itemize}

\begin{longtable}{>{\raggedright\arraybackslash}p{.28\textwidth}>{\raggedright\arraybackslash}p{.64\textwidth}}
\toprule 義務 & 状態と次に必要な根拠 \\
\midrule\endhead
source rank-2 の値型安全性 & 紙上の構成・導出あり。source の通常導出から全称 payload と各 use を生成する。独立査読は未実施。\\
再開の receiver 寿命 & 旧橋は不成立。修正版の局所/全実行型と精密な保存の紙上導出あり。旧効果証拠の自動移行は不可。\\
複数入力の \texttt{'e?} & \(A_m\) の source 解釈を照会中。別 binder/instance の保護を潰すことは既存仕様に反する。\\
捕捉・返却の source 構成 & C-Free/C-Close と定理\ref{thm:capture}の紙上導出あり。実行時の既存証拠を固定する。制約 graph の FreshUse 対応と独立査読は OPEN。\\
compose と non-Var upper-use & 候補 \(\Pi_{\rm flow}\) の alias/actual projection と receiving slot 捕捉を構成。早い receipt へ戻す案を source 判別で棄却。修正版の全期待挙動との一致・採用、native 全経路と主型性は OPEN。\\
一般の annotation/source 生成 & 候補体系の omission seed/upper exposure/closure/target を構成。候補の選択、source の constructor application 一般化を含む受理範囲、全 source expansion は OPEN。\\
source 型健全性 & 固定した \(\Pi_{\rm flow},A\) の体系族で定理\ref{thm:source-safety-family}へ合成。独立査読・機械化は未実施。政策未選択の Yulang 全体の閉包とは異なる。\\
dot selection & 有限証拠からの生成 code は上記定理へ接続。snapshot coherence はこの健全性と独立。resolver 自体の compact/raw・SCC 公開・完全性は OPEN。\\
全表層構文・生成宣言 & 本稿は元の13式・11 pattern まで。parser の全形、control/rule/Yumark/生成 operation 等からの展開対応は別途必要。\\
推論器・native & 主型性、探索完全性/停止性、現在の compiler の全受理、VM/native 対応、原 A/B は OPEN。\\
\bottomrule
\end{longtable}

元の全言語健全性を閉じるには、この表の意味論側の OPEN を解消し、
定理\ref{thm:source-safety-family}の政策/表層範囲を意図した言語と一致させる必要がある。
推論器や resolver の探索完全性を、推論器独立の型健全性の追加前提にはしない。
「独立査読だけが残る」「推論器さえ実装すれば済む」とは結論しない。
今回の成果は source 値型の構成、既存 owner 機械に合わない橋の除去、
直接の compose の origin/target 構成、source からの捕捉生成、
および明記した候補政策族に対する source 型健全性の合成である。
単独著者の紙上導出であり、PDF 組版、構文確認、有限例の実行は数学的検証の代用ではない。

\end{document}
